% Equations du 1er degré dans IRXIR — Mathématiques 3ème — Cours signé OnBuch+
% Programme officiel MINESEC. Compiler avec : tectonic N11-equations-1er-degre-irxir.tex
\newcommand{\DOCMATIERE}{Mathématiques}
\newcommand{\DOCNIVEAU}{3ème}
\newcommand{\DOCMODULE}{Mathématiques – classe de 3ème}
\newcommand{\DOCLECON}{N11}
\newcommand{\DOCTITRE}{Equations du 1er degré dans IRXIR}
% OnBuch+ — préambule « cours signés » ENRICHI (compilé avec tectonic / XeLaTeX)
% Système visuel « Pop » d'OnBuch+ : fond crème (jamais blanc), bordures encre
% épaisses, ombres « dures » décalées, titres Archivo Black en capitales.
% Version MATHÉMATIQUES : graphes (pgfplots/tikz), boîtes pédagogiques
% (définition, propriété, méthode, attention, le savais-tu, exercice résolu,
% exercice, TP, bilan), page de garde et signature OnBuch+ ancrée en bas.
%
% Macros attendues dans main.tex AVANT \input{preamble} :
%   \DOCMATIERE  \DOCNIVEAU  \DOCMODULE  \DOCLECON (numéro)  \DOCTITRE
\documentclass[11pt]{article}

\usepackage{fontspec}
\usepackage[french]{babel}
\usepackage[a4paper,margin=2cm,top=3.4cm,bottom=2.8cm,headheight=1.4cm,footskip=1.3cm]{geometry}
\usepackage{amsmath,amssymb,amsfonts}
\usepackage{mathtools} % \xmapsto, \coloneqq... souvent utilisés par les modèles
\usepackage{cancel} % simplifications barrées
% \abs{z} (module, valeur absolue) et \norm{u} (norme), utilisés par les modèles
\providecommand{\abs}{}\let\abs\relax\DeclarePairedDelimiter{\abs}{\lvert}{\rvert}
\providecommand{\norm}{}\let\norm\relax\DeclarePairedDelimiter{\norm}{\lVert}{\rVert}
\usepackage[table]{xcolor} % [table] : fournit aussi \rowcolors (alternance de lignes)
\usepackage{pagecolor}
\usepackage{tikz}
\usetikzlibrary{arrows.meta,positioning,calc,shapes.geometric,shapes.misc,shapes.arrows,decorations.pathmorphing,decorations.pathreplacing,decorations.markings,patterns,shadows,mindmap,trees,fit,backgrounds,matrix,intersections,angles,quotes}
\usepackage{pgfplots}
\usepackage{pgf-pie} % diagrammes circulaires \pie (statistiques)
\pgfplotsset{compat=1.18}
\usepgfplotslibrary{fillbetween,groupplots} % aires entre courbes (calcul intégral)
\usepackage{enumitem}
\usepackage{fancyhdr}
% [most] charge toutes les bibliothèques tcolorbox (listings, minted, raster,
% documentation...) et gonfle inutilement la mémoire TeX sur un document long
% (~300 boîtes) : on ne charge que ce qui est réellement utilisé.
\usepackage[skins,breakable]{tcolorbox}
\usepackage{graphicx}
\usepackage{colortbl}
\usepackage{booktabs}
\usepackage{tabularx}
\usepackage{diagbox} % case d'en-tête diagonale des tableaux à double entrée
% Colonnes extensibles centrée / gauche / droite (C, L, R), souvent utilisées par les modèles
\newcolumntype{C}{>{\centering\arraybackslash}X}
\newcolumntype{L}{>{\raggedright\arraybackslash}X}
\newcolumntype{R}{>{\raggedleft\arraybackslash}X}
\usepackage{array}
\usepackage{multicol}
\usepackage{siunitx}
% parse-numbers=false : accepte aussi \SI{2,5 \times 10^{-3}}{...}
\sisetup{locale=FR,output-decimal-marker={,},group-minimum-digits=4,parse-numbers=false}
\DeclareSIUnit\ohm{\ensuremath{\symup{\Omega}}} % U+2126 absent de la police
\usepackage{qrcode}
% \degree : commande fréquemment utilisée par les modèles pour un angle ou une
% température en dehors de \SI{}{} (non fournie par les paquets chargés).
\providecommand{\degree}{^{\circ}}
% \qed (fin de démonstration), utilisé par les modèles sans amsthm
\providecommand{\qed}{\hfill$\square$}
% \barre : double barre des valeurs interdites dans un tableau de variations (tkz-tab)
\providecommand{\barre}{\ensuremath{\|}}
% environnement proof (amsthm non chargé) utilisé par les modèles
\ifdefined\proof\else\newenvironment{proof}[1][Démonstration]{\par\noindent\textit{#1.}\ }{\qed\par}\fi
\usepackage{lastpage}
\usepackage{hyperref}
\hypersetup{hidelinks}

% ============================================================
% Palette « Pop » OnBuch+ — mêmes hex que la classe P (plus_ui.dart)
% ============================================================
\definecolor{poporange}{HTML}{F59321}
\definecolor{popdark}{HTML}{E07A0C}
\definecolor{popcream}{HTML}{FFF6E8}
\definecolor{popink}{HTML}{1C1714}
\definecolor{poppurple}{HTML}{7A5AE0}
\definecolor{popgreen}{HTML}{2E9E6A}
\definecolor{poppink}{HTML}{E0457B}
\definecolor{popblue}{HTML}{2F7DE1}
\definecolor{popgold}{HTML}{C98A12}
\definecolor{popmuted}{HTML}{8A7F76}
\definecolor{popline}{HTML}{EADFD2}
\definecolor{popgray}{HTML}{8A7F76}
\colorlet{popgrey}{popgray}
% Alias tolérants (le modèle nomme parfois les couleurs différemment)
\colorlet{popwhite}{white}
\colorlet{popblack}{popink}
\colorlet{popred}{poppink}
\colorlet{popyellow}{popgold}
% Teintes claires pour fonds de tableaux / zones
\colorlet{popblueL}{popblue!10!white}
\colorlet{poporangeL}{poporange!14!white}
\colorlet{popgreenL}{popgreen!12!white}
\colorlet{poppurpleL}{poppurple!12!white}
\colorlet{poppinkL}{poppink!12!white}
\colorlet{popgoldL}{popgold!14!white}

\pagecolor{popcream}

% ============================================================
% Typographie
% ============================================================
\defaultfontfeatures{Ligatures=TeX}
\setmainfont{PlusJakartaSans-Medium.ttf}[Path=../../../fonts/,BoldFont=PlusJakartaSans-Bold.ttf]
% Maths en sans-serif (Fira Math) avec chiffres et lettres droites de Plus
% Jakarta, pour que les formules se fondent dans le texte.
\usepackage[math-style=ISO,bold-style=ISO]{unicode-math}
\setmathfont{FiraMath-Regular.otf}[Path=../../../fonts/]
\setmathfont{PlusJakartaSans-Medium.ttf}[Path=../../../fonts/,range={up/num,up/{Latin,latin},\mathup}]
\usepackage{icomma} % virgule décimale sans espace en mode math
% Caractères mathématiques Unicode absents des polices de texte (Plus Jakarta,
% Archivo Black) : rendus par leur équivalent mathématique, en texte comme en
% formule — sinon ils s'affichent en carré vide (ex. « Arithmétique dans ℤ »).
\usepackage{newunicodechar}
\newunicodechar{ℝ}{\ensuremath{\mathbb{R}}}
\newunicodechar{ℤ}{\ensuremath{\mathbb{Z}}}
\newunicodechar{ℂ}{\ensuremath{\mathbb{C}}}
\newunicodechar{ℕ}{\ensuremath{\mathbb{N}}}
\newunicodechar{ℚ}{\ensuremath{\mathbb{Q}}}
\newunicodechar{𝔻}{\ensuremath{\mathbb{D}}}
\newunicodechar{α}{\ensuremath{\alpha}}
\newunicodechar{β}{\ensuremath{\beta}}
\newunicodechar{γ}{\ensuremath{\gamma}}
\newunicodechar{δ}{\ensuremath{\delta}}
\newunicodechar{ε}{\ensuremath{\varepsilon}}
\newunicodechar{θ}{\ensuremath{\theta}}
\newunicodechar{λ}{\ensuremath{\lambda}}
\newunicodechar{μ}{\ensuremath{\mu}}
\newunicodechar{σ}{\ensuremath{\sigma}}
\newunicodechar{τ}{\ensuremath{\tau}}
\newunicodechar{φ}{\ensuremath{\varphi}}
\newunicodechar{ψ}{\ensuremath{\psi}}
\newunicodechar{ω}{\ensuremath{\omega}}
\newunicodechar{Σ}{\ensuremath{\Sigma}}
% majuscules produites par \MakeUppercase dans les titres (α-aminés -> Α)
\newunicodechar{Α}{\ensuremath{\alpha}}
\newunicodechar{Β}{\ensuremath{\beta}}
\newunicodechar{Γ}{\ensuremath{\Gamma}}
\newunicodechar{Δ}{\ensuremath{\Delta}}
\newunicodechar{Ε}{\ensuremath{\varepsilon}}
\newunicodechar{Θ}{\ensuremath{\Theta}}
\newunicodechar{Λ}{\ensuremath{\Lambda}}
\newunicodechar{Μ}{\ensuremath{\mu}}
\newunicodechar{Τ}{\ensuremath{\tau}}
\newunicodechar{Φ}{\ensuremath{\Phi}}
\newunicodechar{Ψ}{\ensuremath{\Psi}}
\newunicodechar{Ω}{\ensuremath{\Omega}}
\newunicodechar{↦}{\ensuremath{\mapsto}}
\newunicodechar{∈}{\ensuremath{\in}}
\newunicodechar{∉}{\ensuremath{\notin}}
\newunicodechar{∧}{\ensuremath{\wedge}}
\newunicodechar{⇔}{\ensuremath{\Leftrightarrow}}
\newunicodechar{⟺}{\ensuremath{\Longleftrightarrow}}
\newunicodechar{⇒}{\ensuremath{\Rightarrow}}
\newunicodechar{⟹}{\ensuremath{\Longrightarrow}}
\newunicodechar{∘}{\ensuremath{\circ}}
\newunicodechar{∪}{\ensuremath{\cup}}
\newunicodechar{∩}{\ensuremath{\cap}}
\newunicodechar{≡}{\ensuremath{\equiv}}
\newunicodechar{⊂}{\ensuremath{\subset}}
\newunicodechar{⊥}{\ensuremath{\perp}}
\newunicodechar{∃}{\ensuremath{\exists}}
\newunicodechar{∀}{\ensuremath{\forall}}
\newunicodechar{∛}{\ensuremath{\sqrt[3]{\,}}}
\newunicodechar{∜}{\ensuremath{\sqrt[4]{\,}}}
\newunicodechar{⃗}{\ensuremath{{}^{\to}}}
\newunicodechar{ⁿ}{\ifmmode^{n}\else\textsuperscript{\ensuremath{n}}\fi}
\newunicodechar{ˣ}{\ifmmode^{x}\else\textsuperscript{\ensuremath{x}}\fi}
\newunicodechar{ᵇ}{\ifmmode^{b}\else\textsuperscript{\ensuremath{b}}\fi}
\newunicodechar{ʳ}{\ifmmode^{r}\else\textsuperscript{\ensuremath{r}}\fi}
\newunicodechar{ᵘ}{\ifmmode^{u}\else\textsuperscript{\ensuremath{u}}\fi}
\newunicodechar{ʸ}{\ifmmode^{y}\else\textsuperscript{\ensuremath{y}}\fi}
\newunicodechar{ᵃ}{\ifmmode^{a}\else\textsuperscript{\ensuremath{a}}\fi}
\newunicodechar{ᵉ}{\ifmmode^{e}\else\textsuperscript{\ensuremath{e}}\fi}
\newunicodechar{ᵏ}{\ifmmode^{k}\else\textsuperscript{\ensuremath{k}}\fi}
\newunicodechar{ᵖ}{\ifmmode^{p}\else\textsuperscript{\ensuremath{p}}\fi}
\newunicodechar{ᴺ}{\ifmmode^{N}\else\textsuperscript{\ensuremath{N}}\fi}
\newunicodechar{ᶜ}{\ifmmode^{c}\else\textsuperscript{\ensuremath{c}}\fi}
\newunicodechar{ʰ}{\ifmmode^{h}\else\textsuperscript{\ensuremath{h}}\fi}
\newunicodechar{ᵠ}{\ifmmode^{q}\else\textsuperscript{\ensuremath{q}}\fi}
\newunicodechar{ₙ}{\ifmmode_{n}\else\textsubscript{\ensuremath{n}}\fi}
\newunicodechar{ₐ}{\ifmmode_{a}\else\textsubscript{\ensuremath{a}}\fi}
\newunicodechar{ᵢ}{\ifmmode_{i}\else\textsubscript{\ensuremath{i}}\fi}
\newunicodechar{ᵣ}{\ifmmode_{r}\else\textsubscript{\ensuremath{r}}\fi}
\newunicodechar{ₖ}{\ifmmode_{k}\else\textsubscript{\ensuremath{k}}\fi}
\newunicodechar{ᵦ}{\ifmmode_{\beta}\else\textsubscript{\ensuremath{\beta}}\fi}
\newunicodechar{ₓ}{\ifmmode_{x}\else\textsubscript{\ensuremath{x}}\fi}
\newfontfamily\popdisplay{ArchivoBlack-Regular.ttf}[Path=../../../fonts/]
\newfontfamily\poplogo{SpaceGrotesk-Bold.ttf}[Path=../../../fonts/]
\newfontfamily\popsemi{PlusJakartaSans-SemiBold.ttf}[Path=../../../fonts/]
\newfontfamily\popxbold{PlusJakartaSans-ExtraBold.ttf}[Path=../../../fonts/]
\newfontfamily\popxboldls{PlusJakartaSans-ExtraBold.ttf}[Path=../../../fonts/,LetterSpace=3.0]

\setlist[enumerate]{leftmargin=1.7em,itemsep=5pt,topsep=4pt}
\setlist[itemize]{leftmargin=1.7em,itemsep=4pt,topsep=4pt,label={\color{poporange}\textbullet}}
\setlength{\parindent}{0pt}
\setlength{\parskip}{5pt}
\renewcommand{\arraystretch}{1.25}

% pgfplots : style Pop par défaut
\pgfplotsset{
  every axis/.append style={axis line style={popink,line width=1pt},tick style={popink},
    grid style={popline,line width=0.5pt},label style={font=\small},tick label style={font=\footnotesize}},
  popcurve/.style={poporange,line width=1.8pt,no markers,smooth},
}
% Même style disponible dans un \draw TikZ simple (hors pgfplots)
\tikzset{popcurve/.style={poporange,line width=1.8pt}}

% ============================================================
% Pill / squiggle
% ============================================================
\newcommand{\poppill}[3]{%
  \tikz[baseline=-0.55ex]\node[draw=popink,line width=1.2pt,rounded corners=2.6pt,%
    fill=#2,inner xsep=6.5pt,inner ysep=3pt]{%
    \popxboldls\fontsize{7}{7}\selectfont\color{#3}\MakeUppercase{#1}};%
}
\newcommand{\popsquiggle}[1]{%
  \tikz[baseline=-0.4ex]{
    \draw[#1,line width=2.6pt,line cap=round]
      plot [smooth] coordinates {(0,0.09) (0.34,0.01) (0.68,0.21) (1.02,0.01) (1.36,0.10)};
  }%
}
% Mot-clé surligné (marqueur orange sous le texte)
\newcommand{\cle}[1]{{\popxbold\color{popdark}#1}}

% ============================================================
% En-tête / pied
% ============================================================
\pagestyle{fancy}
\fancyhf{}
\renewcommand{\headrulewidth}{0pt}
\renewcommand{\footrulewidth}{0pt}
\lhead{\raisebox{-0.15ex}{\poplogo\fontsize{13}{13}\selectfont\color{popink}OnBuch\color{poporange}+}}
\rhead{\poppill{\DOCMATIERE\ \textbullet\ \DOCNIVEAU\ \textbullet\ Leçon \DOCLECON}{white}{popink}}
\lfoot{\popxbold\fontsize{7.6}{7.6}\selectfont\color{popmuted}\MakeUppercase{Cours signé OnBuch+}\ \textbullet\ {\popsemi\DOCTITRE}}
\rfoot{\tikz[baseline=-0.6ex]\node[rounded corners=8.5pt,fill=popink,minimum height=17pt,minimum width=17pt,inner xsep=5pt,inner sep=0]{\popxbold\fontsize{8}{8}\selectfont\color{popcream}\thepage\,/\,\pageref*{LastPage}};}

% ============================================================
% Titres
% ============================================================
\newcommand{\chaptitre}[2]{%
  \begin{tcolorbox}[enhanced,colback=poporange,colframe=popink,boxrule=2.6pt,arc=11pt,
      drop shadow=popink,left=15pt,right=15pt,top=12pt,bottom=11pt,before skip=3pt,after skip=18pt]
    \poppill{#2}{popink}{popcream}\par\vspace{9pt}
    {\raggedright\hyphenpenalty=10000\exhyphenpenalty=10000\popdisplay\fontsize{22}{24}\selectfont\color{popcream}\MakeUppercase{#1}\par}
  \end{tcolorbox}
}
% Section numérotée : pastille orange avec le numéro + titre Archivo
\newcounter{popsec}
\newcommand{\coursec}[1]{%
  \stepcounter{popsec}\setcounter{popsub}{0}%
  \par\vspace{16pt}\noindent
  \begin{minipage}{\linewidth}
  \tikz[baseline=-0.7ex]\node[draw=popink,line width=1.6pt,fill=poporange,rounded corners=4pt,minimum size=22pt,inner sep=0]{\popdisplay\fontsize{12}{12}\selectfont\color{popcream}\thepopsec};%
  \hspace{8pt}{\popdisplay\fontsize{15}{17}\selectfont\color{popink}\MakeUppercase{#1}}\par
  \vspace{4pt}\noindent{\color{poporange}\rule{36pt}{3.6pt}}
  \end{minipage}\par\nopagebreak\vspace{8pt}
}
\newcounter{popsub}
\newcommand{\courssub}[1]{%
  \stepcounter{popsub}%
  \par\vspace{9pt}\noindent{\popxbold\fontsize{11.5}{13}\selectfont\color{popink}{\color{poporange}\thepopsec.\thepopsub}\hspace{6pt}#1}\par\nopagebreak\vspace{4pt}
}

% ============================================================
% Boîtes pédagogiques — même recette : fond blanc, bordure épaisse colorée,
% ombre dure encre, badge plein. Titre optionnel [#1].
% ============================================================
\tcbset{popbase/.style={breakable,enhanced,colback=white,boxrule=2.3pt,arc=9pt,
  shadow={3pt}{-3pt}{0pt}{popink},left=14pt,right=14pt,top=11pt,bottom=11pt,before skip=11pt,after skip=15pt,
  fontupper=\popsemi\fontsize{10.6}{14.5}\selectfont\color{popink},
  fontlower=\popsemi\fontsize{10.6}{14.5}\selectfont\color{popink}}}
% Coche dessinée (le glyphe ✓ n'existe pas dans les polices du document)
\newcommand{\popcheck}[1]{\tikz[baseline=-0.5ex]\draw[#1,line width=1.6pt,line cap=round,line join=round](-0.13,0.0)--(-0.04,-0.09)--(0.13,0.1);}
\newcommand{\popboxhead}[3]{\poppill{#1}{#2}{white}\ifx\relax#3\relax\else\hspace{6pt}{\popxbold\fontsize{10.6}{12}\selectfont\color{popink}#3}\fi\par\vspace{7pt}}

\newtcolorbox{aretenir}[1][]{popbase,colframe=popgreen,before upper={\popboxhead{À retenir}{popgreen}{#1}}}
\newtcolorbox{exemplebox}[1][]{popbase,colframe=poporange,before upper={\popboxhead{Exemple}{poporange}{#1}}}
\newtcolorbox{definition}[1][]{popbase,colframe=popblue,before upper={\popboxhead{Définition}{popblue}{#1}}}
\newtcolorbox{propriete}[1][]{popbase,colframe=poppurple,before upper={\popboxhead{Propriété}{poppurple}{#1}}}
\newtcolorbox{methode}[1][]{popbase,colframe=popgold,before upper={\popboxhead{Méthode}{popgold}{#1}}}
\newtcolorbox{attention}[1][]{popbase,colframe=poppink,before upper={\popboxhead{Attention}{poppink}{#1}}}
\newtcolorbox{experience}[1][]{popbase,colframe=popblue,colback=popblueL,before upper={\popboxhead{Expérience}{popblue}{#1}}}
% « Le savais-tu ? » : carte violette pleine (contexte, culture, Cameroun)
\newtcolorbox{savaistu}[1][]{popbase,colback=poppurple,colframe=popink,
  fontupper=\popsemi\fontsize{10.4}{14.2}\selectfont\color{white},
  before upper={\poppill{Le savais-tu\,?}{white}{poppurple}\ifx\relax#1\relax\else\hspace{6pt}{\popxbold\color{white}#1}\fi\par\vspace{7pt}}}
% Exercice résolu : énoncé en haut, \tcblower puis la solution sur fond crème
\newcounter{popexo}\newcounter{popexr}
\newtcolorbox{exoresolu}[1][]{popbase,colframe=poporange,colbacklower=popcream,
  segmentation style={popink,line width=1.2pt,dashed},
  before upper={\stepcounter{popexr}\popboxhead{Exercice résolu \thepopexr}{poporange}{#1}},
  before lower={\poppill{Solution}{popink}{popcream}\par\vspace{6pt}}}
% Exercice (non résolu en place) — la correction est dans la partie Corrigés
\newtcolorbox{exercice}[1][]{popbase,colframe=popink,
  before upper={\stepcounter{popexo}\popboxhead{Exercice \thepopexo}{popink}{#1}}}
% Corrigé
\newtcolorbox{corrige}[1][]{popbase,colframe=popgreen,colback=popgreenL,
  before upper={\popboxhead{Corrigé}{popgreen}{#1}}}
% Objectifs / prérequis / bilan
\newtcolorbox{objectifs}{popbase,colframe=popink,colback=poporangeL,
  before upper={\popboxhead{Objectifs de la leçon}{popink}{}}}
\newtcolorbox{prerequis}{popbase,colframe=popink,colback=popblueL,
  before upper={\popboxhead{Prérequis}{popblue}{}}}
\newtcolorbox{bilan}{popbase,colframe=popink,colback=white,boxrule=2.8pt,
  before upper={\popboxhead{Fiche bilan}{poporange}{L'essentiel à connaître pour l'examen}}}
% Figure encadrée avec légende
\newtcolorbox{popfigure}[1][]{enhanced,breakable=false,colback=white,colframe=popline,boxrule=1.4pt,arc=8pt,
  left=10pt,right=10pt,top=10pt,bottom=8pt,before skip=10pt,after skip=12pt,halign=center,
  lower separated=false,halign lower=center,
  fontlower=\popsemi\fontsize{9}{11.5}\selectfont\color{popmuted},
  #1}
\newcommand{\legende}[1]{\par\vspace{6pt}{\popsemi\fontsize{9}{11.5}\selectfont\color{popmuted}#1}}

% ============================================================
% Signature OnBuch+
% ============================================================
\newcommand{\poplogoinline}[1]{{\poplogo\fontsize{#1}{#1}\selectfont\color{popink}OnBuch\color{poporange}+}}
\newcommand{\signatureonbuch}{%
  \begin{tcolorbox}[enhanced,colback=popink,colframe=popink,boxrule=2.2pt,arc=10pt,
      drop shadow=poporange,left=14pt,right=14pt,top=10pt,bottom=10pt,before skip=0pt,after skip=0pt]
    \tikz[baseline=-0.7ex]\node[circle,fill=popgreen,draw=popcream,line width=1.3pt,minimum size=22pt,inner sep=0]{\popcheck{white}};%
    \hspace{9pt}\begin{minipage}[c]{0.62\linewidth}
      {\popxbold\fontsize{12}{13}\selectfont\color{popcream}Signé\ \textbullet\ {\poplogo OnBuch\color{poporange}+}}\par\vspace{2pt}
      {\raggedright\popsemi\fontsize{8}{10.5}\selectfont\color{popline}\MakeUppercase{Équipe pédagogique OnBuch+}\\\MakeUppercase{Conforme au programme officiel MINESEC}\par}
    \end{minipage}\hfill
    {\poplogo\fontsize{22}{22}\selectfont\color{popcream}OnBuch\color{poporange}+}
  \end{tcolorbox}%
}

% ============================================================
% Page de garde
% #1 titre · #2 matière · #3 niveau · #4 module · #5 n° leçon · #6 durée ·
% #7 description / objectifs (texte ou itemize)
% ============================================================
\newcommand{\pagegarde}[7]{%
  \thispagestyle{empty}
  \newgeometry{a4paper,margin=1.7cm,top=1.6cm,bottom=1.4cm}%
  \begin{tikzpicture}[remember picture,overlay]
    \fill[poporange!18] ([xshift=-3.2cm,yshift=-3.6cm]current page.north east) circle (4.6cm);
    \fill[poppurple!14] ([xshift=2.4cm,yshift=7.6cm]current page.south west) circle (3.8cm);
  \end{tikzpicture}%
  {\poplogo\fontsize{30}{30}\selectfont\color{popink}OnBuch\color{poporange}+}\hfill
  \raisebox{0.6ex}{\poppill{Cours signé \textbullet\ Premium}{popink}{popcream}}\par
  \vspace{1pt}\popsquiggle{poppurple}\par
  \vspace{8pt}
  \begin{tcolorbox}[enhanced,colback=poporange,colframe=popink,boxrule=2.8pt,arc=13pt,
      drop shadow=popink,left=18pt,right=18pt,top=12pt,bottom=13pt,after skip=10pt]
    \poppill{#2 \textbullet\ #3}{popink}{popcream}\hspace{5pt}\poppill{#4}{white}{popink}\par\vspace{8pt}
    {\popdisplay\fontsize{14}{14}\selectfont\color{popink}LEÇON #5}\par\vspace{4pt}
    {\raggedright\hyphenpenalty=10000\exhyphenpenalty=10000\popdisplay\fontsize{25}{27}\selectfont\color{popcream}\MakeUppercase{#1}\par}
    \vspace{8pt}
    {\popxbold\fontsize{9.5}{9.5}\selectfont\color{popink}\MakeUppercase{Durée conseillée : #6}}
  \end{tcolorbox}
  \begin{tcolorbox}[enhanced,colback=white,colframe=popink,boxrule=2.2pt,arc=10pt,
      drop shadow=popink,left=16pt,right=16pt,top=10pt,bottom=10pt,after skip=8pt]
    \poppill{Dans cette leçon}{poppurple}{white}\par\vspace{5pt}
    \popsemi\fontsize{9.3}{12.6}\selectfont\color{popink}#7
  \end{tcolorbox}
  \noindent\begin{minipage}[t]{0.487\linewidth}
    \begin{tcolorbox}[enhanced,colback=popgreen,colframe=popink,boxrule=2.2pt,arc=10pt,
        drop shadow=popink,left=11pt,right=11pt,top=10pt,bottom=10pt,height=3.1cm,valign=center]
      \tcbox[colback=white,colframe=popink,boxrule=1.3pt,arc=5pt,boxsep=0pt,nobeforeafter,
        left=3pt,right=3pt,top=3pt,bottom=3pt,valign=center]{\qrcode[height=1.9cm]{https://play.google.com/store/apps/details?id=com.luvvix.onbuch&hl=fr}}%
      \hspace{8pt}\begin{minipage}[c]{0.5\linewidth}
        {\popdisplay\fontsize{11}{12.5}\selectfont\color{white}\MakeUppercase{Télécharge OnBuch}}\par\vspace{4pt}
        {\raggedright\popsemi\fontsize{8.4}{11}\selectfont\color{white}Gratuit \textbullet\ Google Play\\Tuteur IA, annales, concours, résultats\par}
      \end{minipage}
    \end{tcolorbox}
  \end{minipage}\hfill
  \begin{minipage}[t]{0.487\linewidth}
    \begin{tcolorbox}[enhanced,colback=poppurple,colframe=popink,boxrule=2.2pt,arc=10pt,
        drop shadow=popink,left=11pt,right=11pt,top=10pt,bottom=10pt,height=3.1cm,valign=center]
      \tikz[baseline=-0.7ex]\node[circle,fill=white,draw=popink,line width=1.3pt,minimum size=1.6cm,inner sep=0]{\popdisplay\fontsize{24}{24}\selectfont\color{poppurple}+};%
      \hspace{8pt}\begin{minipage}[c]{0.6\linewidth}
        {\popdisplay\fontsize{11}{12.5}\selectfont\color{white}\MakeUppercase{Passe à OnBuch+}}\par\vspace{4pt}
        {\raggedright\popsemi\fontsize{8.4}{11}\selectfont\color{white}Tous les cours signés, Tuteur IA illimité, fiches PDF — onglet OnBuch+ de l'app\par}
      \end{minipage}
    \end{tcolorbox}
  \end{minipage}
  \vspace{8pt}
  \popcontenu
  \vfill
  \signatureonbuch
  \restoregeometry
  \newpage
}

% Rangée « Ce cours contient » (page de garde)
\newcommand{\poptile}[3]{% #1 couleur #2 gros texte #3 légende
  \begin{tcolorbox}[enhanced,colback=#1,colframe=popink,boxrule=2pt,arc=9pt,shadow={2.5pt}{-2.5pt}{0pt}{popink},
      left=6pt,right=6pt,top=9pt,bottom=9pt,halign=center,valign=center,height=1.75cm,width=\linewidth,nobeforeafter]
    {\popdisplay\fontsize{12.5}{13}\selectfont\color{white}\MakeUppercase{#2}}\par\vspace{3pt}
    {\centering\popsemi\fontsize{7.8}{9.5}\selectfont\color{white}#3\par}
  \end{tcolorbox}}
\newcommand{\popcontenu}{%
  \par\noindent{\popxboldls\fontsize{7.5}{7.5}\selectfont\color{popmuted}CE COURS SIGNÉ CONTIENT}\par\vspace{7pt}
  \noindent\begin{minipage}[t]{0.235\linewidth}\poptile{popblue}{Cours}{complet, illustré, pas à pas}\end{minipage}\hfill
  \begin{minipage}[t]{0.235\linewidth}\poptile{poporange}{Méthodes}{et exercices résolus}\end{minipage}\hfill
  \begin{minipage}[t]{0.235\linewidth}\poptile{poppink}{Exercices}{type examen + corrigés}\end{minipage}\hfill
  \begin{minipage}[t]{0.235\linewidth}\poptile{popgreen}{Fiche bilan}{l'essentiel à retenir}\end{minipage}\par}

% Sommaire
\newtcolorbox{sommaire}{enhanced,colback=white,colframe=popink,boxrule=2.2pt,arc=10pt,
  drop shadow=popink,left=18pt,right=18pt,top=13pt,bottom=13pt,before skip=6pt,after skip=20pt,
  before upper={\poppill{Sommaire}{poppurple}{white}\par\vspace{9pt}\popsemi\fontsize{10.6}{14.5}\selectfont\color{popink}}}

% Signature de fin de cours — poussée tout en bas de la dernière page
\newcommand{\findecours}{%
  \par\vfill\nopagebreak
  \begin{center}\popsquiggle{poporange}\end{center}\vspace{4pt}
  \signatureonbuch
}
\AtBeginDocument{\def\llbracket{\mathopen{[\mkern-3.5mu[}}\def\rrbracket{\mathclose{]\mkern-3.5mu]}}} % intervalles d'entiers (glyphe ⟦ absent de la police)
% Glyphes absents de Fira Math : redéfinis à partir de glyphes disponibles
\AtBeginDocument{%
  \def\setminus{\mathbin{\backslash}}\let\smallsetminus\setminus
  \def\vdots{\mathord{\vbox{\baselineskip3.2pt\lineskiplimit0pt\kern4pt\hbox{.}\hbox{.}\hbox{.}}}}%
  \def\ddots{\mathinner{\mkern1mu\raise6.5pt\hbox{.}\mkern2mu\raise3.6pt\hbox{.}\mkern2mu\raise0.7pt\hbox{.}\mkern1mu}}%
  \def\triangle{\mathord{\scalebox{0.9}{$\symup{\Delta}$}}}%
  \def\nabla{\mathord{\raisebox{\depth}{\scalebox{1}[-1]{$\symup{\Delta}$}}}}%
  \def\checkmark{\popcheck{popdark}}%
  \def\star{\mathbin{\ast}}\def\bigstar{\mathord{\ast}}%
  \def\diamond{\mathbin{\scalebox{0.6}{$\Diamond$}}}%
  \def\bigcup{\mathop{\mathchoice{\raisebox{-0.3em}{\scalebox{1.7}{$\cup$}}}{\scalebox{1.25}{$\cup$}}{\cup}{\cup}}\displaylimits}%
  \def\bigcap{\mathop{\mathchoice{\raisebox{-0.3em}{\scalebox{1.7}{$\cap$}}}{\scalebox{1.25}{$\cap$}}{\cap}{\cap}}\displaylimits}%
  \def\lesseqqgtr{\mathrel{\lessgtr}}\def\bigoplus{\mathop{\oplus}\displaylimits}%
}
\providecommand{\ssi}{si et seulement si} % abréviation fréquente
\AtBeginDocument{\providecommand{\wideparen}{\overparen}} % arc de cercle

%% FIGFIX-BEGIN v3 — correctifs globaux des figures (docs/onbuch-generation/tools/figfix)
\usepackage{adjustbox}\usepackage{etoolbox}
% toute figure TikZ/pgfplots plus large que la ligne est réduite à la largeur disponible
\BeforeBeginEnvironment{tikzpicture}{\begin{adjustbox}{max width=\linewidth}}
\AfterEndEnvironment{tikzpicture}{\end{adjustbox}}
% légendes pgfplots lisibles par-dessus les courbes
\pgfplotsset{every axis/.append style={legend style={fill=white,fill opacity=0.92,text opacity=1,draw=black!15,inner sep=3pt}}}
% étiquettes d'arêtes (syntaxe edge["..."]) avec fond blanc
\tikzset{every edge quotes/.append style={fill=white,inner sep=1.2pt}}
% bulles de cartes mentales : texte à la ligne dans la bulle, bulle un peu plus grande
\tikzset{every concept/.append style={text width=2.0cm,align=center,minimum size=2.3cm,inner sep=2pt}}
%% FIGFIX-END
\begin{document}

\pagegarde{Equations du 1er degré dans IRXIR}{Mathématiques}{3ème}{Mathématiques – classe de 3ème}{N11}{2 séances (fiche de progression harmonisée)}{Cette leçon introduit les équations du premier degré dans ℝ×ℝ, leur représentation graphique, les méthodes de résolution de systèmes de deux équations (graphique, substitution, élimination) et l'interprétation du nombre de solutions. Elle met l'accent sur la modélisation de problèmes quotidiens et la rédaction rigoureuse du raisonnement.
\vspace{4pt}
\begin{multicols}{2}\small
\begin{itemize}
\item À la fin de cette leçon, je sais écrire une équation du premier degré à deux inconnues modélisant une situation concrète.
\item Je sais représenter graphiquement une équation du premier degré dans un repère orthonormé.
\item Je sais résoudre un système de deux équations du premier degré par la méthode graphique.
\item Je sais résoudre un système par substitution et par élimination (combinaison linéaire).
\item Je sais discuter le nombre de solutions d'un système (unique, aucune, une infinité) à l'aide du déterminant.
\item Je sais rédiger une démonstration complète et justifier chaque étape de ma résolution.
\end{itemize}\end{multicols}}

\begin{sommaire}
\begin{multicols}{2}
\begin{enumerate}
\item 1. Équation du premier degré à deux inconnues : vocabulaire et représentation
\item 2. Système de deux équations : approche graphique
\item 3. Méthodes algébriques : substitution et élimination (combinaison linéaire)
\item 4. Discussion du nombre de solutions : déterminant et cas particuliers
\item 5. Modélisation et résolution de problèmes concrets
\item Activité d'intégration
\item Méthodes et exercices résolus
\item Exercices
\item Corrigés des exercices
\item Fiche bilan
\end{enumerate}
\end{multicols}
\end{sommaire}

\begin{objectifs}
\begin{itemize}
\item À la fin de cette leçon, je sais écrire une équation du premier degré à deux inconnues modélisant une situation concrète.
\item Je sais représenter graphiquement une équation du premier degré dans un repère orthonormé.
\item Je sais résoudre un système de deux équations du premier degré par la méthode graphique.
\item Je sais résoudre un système par substitution et par élimination (combinaison linéaire).
\item Je sais discuter le nombre de solutions d'un système (unique, aucune, une infinité) à l'aide du déterminant.
\item Je sais rédiger une démonstration complète et justifier chaque étape de ma résolution.
\item Je sais appliquer ces outils à des problèmes de la vie courante (budget, mélange, géométrie).
\item Je sais identifier et éviter les erreurs fréquentes (mauvaise manipulation des signes, confusion entre équation et système).
\end{itemize}
\end{objectifs}

\begin{prerequis}
\begin{itemize}
\item Résolution d'équations du premier degré à une inconnue (classe de 4ème).
\item Repérage dans le plan et tracé de droites (abscisse, ordonnée, pente).
\item Notion de couple ordonné (x, y) et lecture de coordonnées sur un graphique.
\item Calculs avec les nombres relatifs et les fractions.
\item Vocabulaire de base : variable, coefficient, terme constant, solution.
\end{itemize}
\end{prerequis}

\newpage

\coursec{Équation du premier degré à deux inconnues : vocabulaire et représentation}

\courssub{Une situation de la vie courante pour ouvrir la leçon}

Maman Ngo prépare la fête de fin d'année de l'école de son quartier à Yaoundé. Elle dispose de \SI{1500}{FCFA} et souhaite acheter des beignets et des boissons. Un beignet coûte \SI{150}{FCFA} et une boisson coûte \SI{200}{FCFA}. Elle veut dépenser \emph{exactement} son budget, sans reste ni dépassement.

Notons $x$ le nombre de beignets et $y$ le nombre de boissons achetés. La dépense totale s'écrit alors :
\[ 150x + 200y = 1500. \]

Voilà une égalité qui contient \textbf{deux} nombres inconnus à la fois. Contrairement aux équations à une seule inconnue que tu connais depuis la classe de 4ème (comme $3x+5=20$), il n'y a plus une seule valeur à trouver, mais des \textbf{couples} de valeurs $(x\,;\,y)$ qui rendent l'égalité vraie. Par exemple, le couple $(2\,;\,6)$ convient : $150\times 2 + 200\times 6 = 300 + 1200 = 1500$. Mais le couple $(6\,;\,3)$ convient aussi : $150\times 6 + 200\times 3 = 900 + 600 = 1500$. Il y a donc \emph{plusieurs} solutions !

\textbf{Problématique de la leçon :} qu'est-ce qu'une équation du premier degré à deux inconnues ? Comment décrire \emph{toutes} ses solutions ? Et pourquoi ces solutions forment-elles une droite dans un repère ?

\courssub{Définition et vocabulaire}

\begin{definition}[Équation du premier degré à deux inconnues]
Une \cle{équation du premier degré à deux inconnues} $x$ et $y$ est une équation qui peut s'écrire sous la forme
\[ ax + by = c \]
où $a$, $b$ et $c$ sont des nombres réels donnés, avec $(a\,;\,b) \neq (0\,;\,0)$, c'est-à-dire que $a$ et $b$ ne sont \textbf{pas tous les deux nuls}.

Une \cle{solution} de cette équation est un couple $(x\,;\,y)$ de nombres réels pour lequel l'égalité $ax+by=c$ est vraie.
\end{definition}

\begin{exemplebox}[Reconnaître une équation du premier degré à deux inconnues]
\begin{itemize}
\item $3x - 2y = 6$ est une équation du premier degré à deux inconnues, avec $a=3$, $b=-2$ et $c=6$.
\item $150x + 200y = 1500$ (le problème de Maman Ngo) en est une aussi.
\item $x + y = 10$ en est une, avec $a=1$, $b=1$, $c=10$.
\item $x^2 + y = 3$ n'en est \textbf{pas} une, à cause du carré sur $x$ (ce n'est plus le premier degré).
\item $xy = 4$ n'en est \textbf{pas} une non plus : le produit des deux inconnues n'est pas autorisé.
\end{itemize}
\end{exemplebox}

\begin{attention}[Pourquoi exclure le cas $a = b = 0$ ?]
Si $a=0$ et $b=0$, l'équation devient $0 = c$. Si $c \neq 0$, aucun couple ne convient ; si $c=0$, tous les couples conviennent. Ce cas ne présente aucun intérêt : c'est pourquoi on exige $(a\,;\,b) \neq (0\,;\,0)$ dans la définition. Retiens aussi qu'une solution est un \textbf{couple ordonné} : le couple $(2\,;\,6)$ n'est pas le même que $(6\,;\,2)$, et il se peut que l'un soit solution et pas l'autre !
\end{attention}

\begin{exoresolu}[Vérifier qu'un couple est solution]
\textbf{Énoncé.} On considère l'équation $3x - 2y = 6$. Les couples $(4\,;\,3)$ et $(1\,;\,2)$ sont-ils solutions ?
\tcblower
\textbf{Solution.} Pour le couple $(4\,;\,3)$, on remplace $x$ par $4$ et $y$ par $3$ :
\[ 3\times 4 - 2\times 3 = 12 - 6 = 6. \]
L'égalité est vraie : $(4\,;\,3)$ \textbf{est solution}.

Pour le couple $(1\,;\,2)$ :
\[ 3\times 1 - 2\times 2 = 3 - 4 = -1 \neq 6. \]
L'égalité est fausse : $(1\,;\,2)$ \textbf{n'est pas solution}.
\end{exoresolu}

\courssub{L'ensemble des solutions : une droite}

Reprenons l'équation $3x - 2y = 6$. On a vu que $(4\,;\,3)$ est solution. On peut en trouver d'autres : $(0\,;\,-3)$, $(2\,;\,0)$, $(6\,;\,6)$... En fait, on peut choisir $x$ librement et calculer le $y$ correspondant : il y a donc une \textbf{infinité} de solutions. Plaçons ces couples dans un repère : surprise, tous les points obtenus sont alignés !

\begin{propriete}[Ensemble des solutions (admise)]
Dans un repère du plan, l'ensemble des points $M(x\,;\,y)$ dont les coordonnées sont solutions de l'équation $ax + by = c$ (avec $(a\,;\,b)\neq(0\,;\,0)$) est une \cle{droite}. On dit que $ax+by=c$ est une \cle{équation de la droite}.
\end{propriete}

\begin{attention}[Équation ou droite ? Ne pas confondre !]
L'\textbf{équation} $ax+by=c$ est une \emph{relation} entre deux nombres $x$ et $y$. L'\textbf{ensemble des solutions} est un objet \emph{géométrique} : une droite du plan. Dire « la solution est la droite » est un abus de langage ; la formulation correcte est : « l'ensemble des couples solutions est représenté par une droite ». Au BEPC, soigne ce vocabulaire : un couple est solution, un point appartient à la droite.
\end{attention}

\courssub{Tracer la droite : deux points suffisent}

Pour tracer une droite, il suffit d'en connaître deux points. L'astuce consiste à choisir les points les plus faciles à calculer : celui où la droite coupe l'axe des ordonnées (on pose $x=0$) et celui où elle coupe l'axe des abscisses (on pose $y=0$).

\begin{methode}[Tracer la droite d'équation $ax+by=c$]
\begin{enumerate}
\item On pose $x=0$ dans l'équation et on résout $by=c$ : on obtient le point $(0\,;\,\tfrac{c}{b})$ si $b\neq 0$ (intersection avec l'axe des ordonnées).
\item On pose $y=0$ et on résout $ax=c$ : on obtient le point $(\tfrac{c}{a}\,;\,0)$ si $a\neq 0$ (intersection avec l'axe des abscisses).
\item On place ces deux points dans un repère et on trace la droite qui les relie, à la règle.
\item Vérification : on calcule un troisième point et on contrôle qu'il est bien sur la droite tracée.
\end{enumerate}
\end{methode}

\begin{exemplebox}[Tracer la droite d'équation $3x - 2y = 6$]
\begin{itemize}
\item Si $x=0$ : $-2y = 6$, donc $y = -3$. Premier point : $(0\,;\,-3)$.
\item Si $y=0$ : $3x = 6$, donc $x = 2$. Deuxième point : $(2\,;\,0)$.
\end{itemize}
On place ces deux points, on trace la droite. Vérification avec $x=4$ : $12-2y=6$ donne $y=3$ ; le point $(4\,;\,3)$ est bien aligné avec les deux autres.
\end{exemplebox}

\begin{exemplebox}[La droite de Maman Ngo]
Pour l'équation $150x + 200y = 1500$, on peut simplifier par $50$ : $3x + 4y = 30$.
\begin{itemize}
\item Si $x=0$ : $4y=30$, donc $y = 7,5$. Point $(0\,;\,7,5)$.
\item Si $y=0$ : $3x=30$, donc $x=10$. Point $(10\,;\,0)$.
\end{itemize}
Tous les points de cette droite représentent les budgets possibles. Mais attention au contexte : Maman Ngo ne peut acheter qu'un nombre \emph{entier} de beignets et de boissons, donc seuls les couples d'entiers positifs de la droite ont un sens concret, comme $(2\,;\,6)$ ou $(6\,;\,3)$.
\end{exemplebox}

\begin{popfigure}
\begin{center}
\begin{tikzpicture}[scale=0.8]
\draw[->] (-0.5,0)--(4.5,0) node[right]{$x$};
\draw[->] (0,-3.5)--(0,3.5) node[above]{$y$};
\draw[thick,popblue] (0,-3)--(2,0) node[midway,above right]{$3x-2y=6$};
\fill (0,-3) circle(2pt) node[left]{$(0\,;\,-3)$};
\fill (2,0) circle(2pt) node[below]{$(2\,;\,0)$};
\fill[poporange] (4,3) circle(2pt) node[above right]{$(4\,;\,3)$};
\end{tikzpicture}
\end{center}
\legende{La droite d'équation $3x-2y=6$ passe par $(0\,;\,-3)$ et $(2\,;\,0)$. Le point $(4\,;\,3)$ est aussi solution : il est sur la droite.}
\end{popfigure}

\courssub{L'équation réduite : $y = mx + p$}

Lorsque $b \neq 0$, on peut isoler $y$ dans l'équation $ax+by=c$ :
\[ by = -ax + c \quad\Longrightarrow\quad y = -\frac{a}{b}\,x + \frac{c}{b}. \]
En posant $m = -\dfrac{a}{b}$ et $p = \dfrac{c}{b}$, on obtient la forme dite \emph{réduite}.

\begin{definition}[Équation réduite d'une droite]
Toute droite non verticale (c'est-à-dire avec $b\neq 0$) admet une \cle{équation réduite} de la forme
\[ y = mx + p \]
où $m$ est le \cle{coefficient directeur} (ou pente) de la droite et $p$ est son \cle{ordonnée à l'origine} : la droite coupe l'axe des ordonnées au point $(0\,;\,p)$.
\end{definition}

\begin{exemplebox}[Passer à la forme réduite]
Reprenons $3x - 2y = 6$. On isole $y$ :
\begin{align*}
-2y &= -3x + 6 \\
y &= \frac{3}{2}x - 3.
\end{align*}
La pente est $m = \dfrac{3}{2}$ et l'ordonnée à l'origine est $p = -3$. On retrouve bien le point $(0\,;\,-3)$ calculé plus haut !

La pente se lit géométriquement : quand $x$ augmente de $2$, $y$ augmente de $3$. C'est exactement ce qu'on observe en passant de $(0\,;\,-3)$ à $(2\,;\,0)$.
\end{exemplebox}

\begin{attention}[Cas des droites verticales et horizontales]
\begin{itemize}
\item Si $b = 0$ (et $a\neq 0$), l'équation devient $ax=c$, soit $x = \dfrac{c}{a}$ : c'est une droite \textbf{verticale}. Elle n'a \emph{pas} de forme réduite $y=mx+p$.
\item Si $a = 0$ (et $b\neq 0$), l'équation devient $y = \dfrac{c}{b}$ : c'est une droite \textbf{horizontale}, de pente $m=0$.
\end{itemize}
Ne cherche donc jamais à écrire $y=mx+p$ pour une droite verticale comme $x=3$ : c'est impossible !
\end{attention}

\courssub{Lecture graphique d'une solution}

Une fois la droite tracée, on peut lire graphiquement des solutions :
\begin{itemize}
\item pour trouver $y$ connaissant $x$ : on part de la valeur de $x$ sur l'axe des abscisses, on monte (ou descend) verticalement jusqu'à la droite, puis on lit l'ordonnée sur l'axe des $y$ ;
\item pour trouver $x$ connaissant $y$ : on fait le trajet inverse, horizontalement puis verticalement.
\end{itemize}
La lecture graphique donne souvent une valeur \emph{approchée} ; le calcul algébrique, lui, donne la valeur \emph{exacte}. Les deux se complètent : le graphique sert à vérifier la cohérence du calcul.

\begin{exercice}[À toi de jouer]
On considère l'équation $2x + 3y = 12$.
\begin{enumerate}
\item Vérifier que le couple $(3\,;\,2)$ est solution.
\item Déterminer deux points de la droite associée, puis la tracer dans un repère.
\item Écrire l'équation réduite de cette droite et donner sa pente.
\item Lire graphiquement la valeur de $y$ lorsque $x = 1,5$, puis vérifier par le calcul.
\end{enumerate}
\end{exercice}

\begin{corrige}[Exercice 1]
\begin{enumerate}
\item $2\times 3 + 3\times 2 = 6 + 6 = 12$ : le couple $(3\,;\,2)$ est bien solution.
\item Si $x=0$ : $3y=12$ donc $y=4$, point $(0\,;\,4)$. Si $y=0$ : $2x=12$ donc $x=6$, point $(6\,;\,0)$. On trace la droite passant par ces deux points.
\item $3y = -2x + 12$ donc $y = -\dfrac{2}{3}x + 4$. La pente est $m=-\dfrac{2}{3}$ et l'ordonnée à l'origine est $p=4$.
\item Graphiquement, pour $x=1,5$ on lit $y=3$. Par le calcul : $y = -\dfrac{2}{3}\times 1,5 + 4 = -1 + 4 = 3$. La lecture est confirmée.
\end{enumerate}
\end{corrige}

\begin{aretenir}[L'essentiel de la section]
\begin{itemize}
\item Une équation du premier degré à deux inconnues s'écrit $ax+by=c$ avec $(a\,;\,b)\neq(0\,;\,0)$ ; ses solutions sont des \textbf{couples} $(x\,;\,y)$.
\item L'ensemble des solutions est représenté par une \textbf{droite} dans un repère.
\item Pour tracer cette droite, on calcule deux points, en posant successivement $x=0$ puis $y=0$.
\item Si $b\neq 0$, la forme réduite $y=mx+p$ donne la pente $m=-\dfrac{a}{b}$ et l'ordonnée à l'origine $p=\dfrac{c}{b}$.
\item Une droite verticale ($x=k$) n'a pas de forme réduite.
\end{itemize}
\end{aretenir}

\begin{savaistu}[Pourquoi « réduite » ?]
On dit que $y=mx+p$ est la forme \emph{réduite} parce que $y$ est « réduit » à être seul dans son membre : on a isolé l'inconnue. Cette écriture est très appréciée des économistes et des commerçants : dans $y = 150x + 500$, on lit immédiatement un coût fixe de \SI{500}{FCFA} (location du stand au marché) et un coût de \SI{150}{FCFA} par article vendu. La pente, c'est le prix unitaire ! Remarque aussi que toute droite non verticale, même tracée au hasard sur une feuille, peut s'écrire sous la forme $y=mx+p$ : il suffit de lire où elle coupe l'axe des ordonnées et de mesurer son inclinaison.
\end{savaistu}

\coursec{Système de deux équations : approche graphique}

Dans la section précédente, nous avons découvert qu'une équation du premier degré à deux inconnues, comme $x + y = 3$, possède une \emph{infinité} de couples solutions, et que l'ensemble de ces couples se représente par une \cle{droite} dans le plan muni d'un repère. Une question naturelle se pose alors : que se passe-t-il si l'on impose \emph{deux conditions à la fois} ? Par exemple, un commerçant de Mokolo vend des cahiers et des stylos : au total 3 articles ont été achetés pour une dépense donnée. Chercher deux nombres vérifiant \emph{simultanément} deux équations, c'est exactement résoudre un \cle{système}. Et géométriquement, cela revient à chercher l'intersection de deux droites. C'est cette approche graphique, intuitive et visuelle, que nous allons explorer dans cette section.

\courssub{Définition d'un système de deux équations}

\begin{definition}[Système de deux équations du premier degré à deux inconnues]
Un \cle{système de deux équations du premier degré à deux inconnues} $x$ et $y$ est la donnée de deux équations du premier degré que l'on cherche à satisfaire \textbf{en même temps}. Il s'écrit avec une accolade :
\[
\begin{cases} ax + by = c \\ a'x + b'y = c' \end{cases}
\]
où $a, b, c, a', b', c'$ sont des nombres réels donnés (avec $a$ et $b$ non tous les deux nuls, de même pour $a'$ et $b'$).

Une \cle{solution} du système est un couple $(x\,;y)$ de réels qui vérifie \textbf{les deux équations simultanément}. \emph{Résoudre} le système, c'est déterminer \textbf{tous} les couples solutions.
\end{definition}

\begin{attention}[L'accolade signifie « et »]
L'accolade est essentielle : elle signifie que le couple cherché doit vérifier la première équation \textbf{et} la seconde. Un couple qui ne vérifie qu'une seule des deux équations n'est \textbf{pas} une solution du système. Par exemple, pour $\begin{cases} x + y = 3 \\ 2x + y = 3 \end{cases}$, le couple $(3\,;0)$ vérifie la première équation ($3+0=3$) mais pas la seconde ($2\times 3 + 0 = 6 \neq 3$) : ce n'est pas une solution du système. Le couple $(0\,;3)$ vérifie les deux : $0+3=3$ et $2\times 0+3=3$, c'est donc la solution.
\end{attention}

\courssub{Interprétation géométrique : l'intersection de deux droites}

Voici l'idée centrale de cette section. Nous avons vu dans la section 1 que :
\begin{itemize}
\item l'ensemble des couples $(x\,;y)$ vérifiant $ax + by = c$ est une droite $(d_1)$ ;
\item l'ensemble des couples vérifiant $a'x + b'y = c'$ est une droite $(d_2)$.
\end{itemize}
Un couple solution du système doit appartenir aux \emph{deux} ensembles à la fois : c'est donc un point situé \textbf{à la fois sur $(d_1)$ et sur $(d_2)$}, c'est-à-dire un \cle{point d'intersection} des deux droites.

\begin{aretenir}[Lien fondamental système $\leftrightarrow$ intersection]
Résoudre graphiquement le système $\begin{cases} ax + by = c \\ a'x + b'y = c' \end{cases}$ revient à :
\begin{enumerate}
\item tracer la droite $(d_1)$ d'équation $ax + by = c$ ;
\item tracer la droite $(d_2)$ d'équation $a'x + b'y = c'$ ;
\item lire les coordonnées du (ou des) point(s) d'intersection.
\end{enumerate}
\end{aretenir}

\begin{methode}[Résolution graphique d'un système, pas à pas]
\begin{enumerate}
\item Pour chaque équation, trouver \textbf{deux points} de la droite (on choisit souvent $x = 0$ puis $y = 0$, ce qui donne les intersections avec les axes).
\item Placer ces points dans un repère et tracer soigneusement les deux droites à la règle.
\item Repérer le point d'intersection éventuel et lire ses coordonnées $(x_0\,;y_0)$ sur les axes.
\item \textbf{Vérifier} en remplaçant $x$ par $x_0$ et $y$ par $y_0$ dans les \emph{deux} équations : les deux égalités doivent être vraies.
\end{enumerate}
\end{methode}

\begin{exemplebox}[Un premier système résolu graphiquement]
Résolvons graphiquement le système $\begin{cases} x + y = 3 \\ 2x + y = 3 \end{cases}$.

\textbf{Droite $(d_1)$ : $x + y = 3$.} Si $x = 0$, alors $y = 3$ ; si $y = 0$, alors $x = 3$. La droite passe par $A(0\,;3)$ et $B(3\,;0)$.

\textbf{Droite $(d_2)$ : $2x + y = 3$.} Si $x = 0$, alors $y = 3$ ; la droite passe par $A(0\,;3)$. Si $y = 0$, alors $2x = 3$ donc $x = 1,5$. La droite passe par $C(1,5\,;0)$.

Les deux droites se coupent en $A(0\,;3)$. Vérification : $0 + 3 = 3$ \checkmark{} et $2\times 0 + 3 = 3$ \checkmark. L'unique solution du système est le couple $(0\,;3)$.
\end{exemplebox}

\begin{popfigure}
\begin{center}
\begin{tikzpicture}[scale=0.8]
% Axes
\draw[->] (-0.5,0) -- (4,0) node[right] {$x$};
\draw[->] (0,-0.5) -- (0,4) node[above] {$y$};
% Graduations
\foreach \x in {1,2,3} \draw (\x,0.05) -- (\x,-0.05) node[below] {\small $\x$};
\foreach \y in {1,2,3} \draw (0.05,\y) -- (-0.05,\y) node[left] {\small $\y$};
% Droite d1 : x+y=3 -> (0,3) to (3,0) extended
\draw[thick, popblue] (0,3) -- (3.5,-0.5) node[below right, pos=0.8] {$x+y=3$};
% Droite d2 : 2x+y=3 -> (0,3) to (1.5,0) extended
\draw[thick, poporange] (0,3) -- (2.5,-2) node[below right, pos=0.7] {$2x+y=3$};
% Point d'intersection
\fill[popink] (0,3) circle(3pt) node[above left] {$(0\,;3)$};
\end{tikzpicture}
\end{center}
\legende{Les droites d'équations $x+y=3$ (bleue) et $2x+y=3$ (orange) se coupent au point $(0\,;3)$ : c'est l'unique solution du système.}
\end{popfigure}

\begin{attention}[Toujours vérifier le point lu]
Sur la figure ci-dessus, on lit le point d'intersection $(0\,;3)$. Cette lecture doit \emph{toujours} être confirmée par le calcul : $0 + 3 = 3$ \checkmark{} et $2\times 0 + 3 = 3$ \checkmark. Moralité : \textbf{ne jamais faire confiance aveuglément à un dessin} ; la vérification algébrique est obligatoire. C'est précisément ce que montre cet exemple : une figure mal graduée ou approximative peut induire en erreur.
\end{attention}

\courssub{Les trois cas possibles}

Deux droites dans le plan ne peuvent se présenter que de trois façons. Ce fait géométrique bien connu depuis la classe de 4ème se traduit directement en termes de nombre de solutions du système.

\begin{propriete}[Nombre de solutions d'un système de deux équations à deux inconnues]
Un système de deux équations du premier degré à deux inconnues possède :
\begin{itemize}
\item \textbf{un unique couple solution} lorsque les deux droites sont \cle{sécantes} (elles se coupent en un seul point) ;
\item \textbf{aucune solution} lorsque les deux droites sont \cle{strictement parallèles} (parallèles et distinctes : elles ne se rencontrent jamais) ;
\item \textbf{une infinité de solutions} lorsque les deux droites sont \cle{confondues} (tous les points de la droite commune sont solutions).
\end{itemize}
\end{propriete}

\begin{exemplebox}[Illustration des trois cas]
\textbf{Cas 1 — droites sécantes.} $\begin{cases} x + y = 3 \\ x - y = 1 \end{cases}$ : les droites se coupent en $(2\,;1)$. Vérification : $2+1=3$ et $2-1=1$. Une seule solution.

\textbf{Cas 2 — droites strictement parallèles.} $\begin{cases} x + y = 3 \\ x + y = 5 \end{cases}$ : les deux droites ont la même direction (même coefficient directeur $-1$) mais passent par des points différents ($(0\,;3)$ et $(0\,;5)$). Elles ne se rencontrent jamais : aucun couple ne peut vérifier à la fois $x+y=3$ et $x+y=5$. Aucune solution.

\textbf{Cas 3 — droites confondues.} $\begin{cases} x + y = 3 \\ 2x + 2y = 6 \end{cases}$ : la seconde équation est simplement la première multipliée par $2$. Les deux équations décrivent la \emph{même} droite. Tous les couples $(x\,;3-x)$ sont solutions : il y en a une infinité.
\end{exemplebox}

\begin{popfigure}
\begin{center}
\begin{tikzpicture}[scale=0.6]
% Cas 1 : sécantes
\begin{scope}
\draw[->] (-0.5,0)--(4,0) node[right]{$x$};
\draw[->] (0,-1.5)--(0,4) node[above]{$y$};
\foreach \x in {1,2,3} \draw (\x,0.05)--(\x,-0.05) node[below]{\tiny $\x$};
\foreach \y in {-1,1,2,3} \draw (0.05,\y)--(-0.05,\y) node[left]{\tiny $\y$};
\draw[thick,popblue] (0,3)--(3.5,-0.5);
\draw[thick,poporange] (0,-1)--(3.5,2.5);
\fill[popink] (2,1) circle(2.5pt);
\node[below] at (1.75,-1.3) {\small Sécantes : 1 solution};
\end{scope}
% Cas 2 : parallèles
\begin{scope}[xshift=6.5cm]
\draw[->] (-0.5,0)--(4,0) node[right]{$x$};
\draw[->] (0,-0.5)--(0,4) node[above]{$y$};
\foreach \x in {1,2,3} \draw (\x,0.05)--(\x,-0.05) node[below]{\tiny $\x$};
\foreach \y in {1,2,3} \draw (0.05,\y)--(-0.05,\y) node[left]{\tiny $\y$};
\draw[thick,popblue] (0,3)--(3.5,-0.5);
\draw[thick,poporange] (0.5,4)--(4,0.5);
\node[below] at (1.75,-0.8) {\small Parallèles : 0 solution};
\end{scope}
% Cas 3 : confondues
\begin{scope}[xshift=13cm]
\draw[->] (-0.5,0)--(4,0) node[right]{$x$};
\draw[->] (0,-0.5)--(0,4) node[above]{$y$};
\foreach \x in {1,2,3} \draw (\x,0.05)--(\x,-0.05) node[below]{\tiny $\x$};
\foreach \y in {1,2,3} \draw (0.05,\y)--(-0.05,\y) node[left]{\tiny $\y$};
\draw[thick,popblue] (0,3)--(3.5,-0.5);
\draw[thick,poporange,dashed] (0.1,2.9)--(3.6,-0.6);
\node[below] at (1.75,-0.8) {\small Confondues : infinité};
\end{scope}
\end{tikzpicture}
\end{center}
\legende{Les trois configurations possibles pour deux droites du plan, et le nombre de solutions correspondant du système.}
\end{popfigure}

\begin{attention}[Reconnaître les cas sans dessiner]
On peut souvent deviner le cas sans tracer :
\begin{itemize}
\item si une équation est un \textbf{multiple exact} de l'autre (coefficients et second membre multipliés par le même nombre non nul), les droites sont \textbf{confondues} ;
\item si les coefficients de $x$ et $y$ sont proportionnels \textbf{mais pas} les seconds membres (par exemple $x+y=3$ et $x+y=5$, ou $2x+2y=6$ et $2x+2y=10$), les droites sont \textbf{strictement parallèles} ;
\item sinon, les droites sont \textbf{sécantes} et le système a une unique solution.
\end{itemize}
Nous étudierons ce critère de manière algébrique (avec le déterminant) dans la section 4.
\end{attention}

\courssub{Lecture approximative sur un graphique}

En pratique, le point d'intersection ne tombe pas toujours sur des coordonnées entières. La méthode graphique permet alors une \cle{lecture approximative}.

\begin{methode}[Lire une solution approchée]
\begin{enumerate}
\item Choisir une échelle adaptée (par exemple 1 cm pour 1 unité, ou 1 cm pour 0,5 unité si l'intersection est « entre les graduations »).
\item Tracer les droites avec soin, au crayon bien taillé, en prolongeant suffisamment.
\item Projeter le point d'intersection sur chaque axe (pointillés) et lire les valeurs approchées.
\item Donner le résultat avec une mention d'approximation : $(x\,;y) \approx (2,3\,;1,7)$, puis vérifier par substitution que les égalités sont « presque » vraies.
\end{enumerate}
Avec un logiciel de géométrie dynamique (GeoGebra, par exemple, disponible sur les smartphones utilisés dans beaucoup d'établissements de Yaoundé ou de Douala), on peut tracer les deux droites et faire afficher les coordonnées exactes du point d'intersection : c'est un excellent moyen de \emph{contrôler} un résultat, mais pas de le \emph{prouver}.
\end{methode}

\begin{exoresolu}[Résolution graphique complète d'un système]
Résoudre graphiquement le système $\begin{cases} x + 2y = 5 \\ 3x - y = 4 \end{cases}$, puis vérifier la solution par le calcul.
\tcblower
\textbf{Étape 1 : points de la droite $(d_1)$ d'équation $x + 2y = 5$.}

Si $x = 0$ : $2y = 5$ donc $y = 2,5$ ; point $A(0\,;2,5)$.
Si $y = 0$ : $x = 5$ ; point $B(5\,;0)$.

\textbf{Étape 2 : points de la droite $(d_2)$ d'équation $3x - y = 4$.}

Si $x = 0$ : $-y = 4$ donc $y = -4$ ; point $C(0\,;-4)$.
Si $y = 0$ : $3x = 4$ donc $x = \dfrac{4}{3} \approx 1,33$ ; point $D\left(\dfrac{4}{3}\,;0\right)$.
Un troisième point plus pratique : si $x = 2$, alors $y = 3\times 2 - 4 = 2$ ; point $E(2\,;2)$.

\textbf{Étape 3 : tracé et lecture.} On trace les deux droites dans un repère ; elles se coupent. La lecture du point d'intersection donne des coordonnées proches de $(1,9\,;1,6)$.

\textbf{Étape 4 : vérification algébrique et recherche de la valeur exacte.}

Le couple approché $(1,9\,;1,6)$ donne : $1,9 + 2\times 1,6 = 5,1 \approx 5$ et $3\times 1,9 - 1,6 = 4,1 \approx 4$. La lecture graphique suggère une solution non entière. Résolvons exactement par substitution (méthode de la section 3) : de la seconde équation $y = 3x - 4$. En remplaçant dans la première : $x + 2(3x - 4) = 5 \iff 7x - 8 = 5 \iff 7x = 13 \iff x = \dfrac{13}{7}$. Alors $y = 3\times\dfrac{13}{7} - 4 = \dfrac{39-28}{7} = \dfrac{11}{7}$.

Vérification exacte : $\dfrac{13}{7} + 2\times\dfrac{11}{7} = \dfrac{35}{7} = 5$ \checkmark{} et $3\times\dfrac{13}{7} - \dfrac{11}{7} = \dfrac{28}{7} = 4$ \checkmark.

\textbf{Conclusion.} Le système admet un unique couple solution : $\left(\dfrac{13}{7}\,;\dfrac{11}{7}\right) \approx (1,86\,;1,57)$. Cet exemple montre à quel point la lecture graphique peut être trompeuse quand la solution n'est pas entière : seule la vérification algébrique (ou une méthode de calcul, voir section 3) donne la valeur exacte.
\end{exoresolu}

\begin{popfigure}
\begin{center}
\begin{tikzpicture}[scale=0.7]
% Axes
\draw[->] (-0.5,0) -- (5.5,0) node[right] {$x$};
\draw[->] (0,-1.5) -- (0,3.5) node[above] {$y$};
% Graduations
\foreach \x in {1,2,3,4,5} \draw (\x,0.05) -- (\x,-0.05) node[below] {\small $\x$};
\foreach \y in {-1,1,2,3} \draw (0.05,\y) -- (-0.05,\y) node[left] {\small $\y$};
% Droite d1 : x+2y=5 -> (0,2.5) to (5,0)
\draw[thick, popblue] (0,2.5) -- (5.5,-0.25) node[below right, pos=0.85] {$x+2y=5$};
% Droite d2 : 3x-y=4 -> y=3x-4. Points (1.33,0), (2,2), (3,5)
\draw[thick, poporange] (1.33,0) -- (2.5,3.5) node[above right, pos=0.7] {$3x-y=4$};
% Intersection (13/7, 11/7) approx (1.857, 1.571)
\fill[popink] (1.857, 1.571) circle(3pt) node[above left] {$\left(\frac{13}{7},\frac{11}{7}\right)$};
% Projections
\draw[dashed, popmuted] (1.857, 1.571) -- (1.857, 0) node[below] {\small $1,86$};
\draw[dashed, popmuted] (1.857, 1.571) -- (0, 1.571) node[left] {\small $1,57$};
\end{tikzpicture}
\end{center}
\legende{Résolution graphique de $\begin{cases} x+2y=5 \\ 3x-y=4 \end{cases}$ : intersection en $\left(\frac{13}{7},\frac{11}{7}\right) \approx (1,86\,;1,57)$.}
\end{popfigure}

\courssub{Limites de la méthode graphique}

La méthode graphique est précieuse pour \emph{comprendre} ce qu'est un système et \emph{visualiser} le nombre de solutions. Mais elle souffre de limites sérieuses qu'il faut connaître.

\begin{attention}[Pourquoi le graphique ne suffit pas]
\begin{itemize}
\item \textbf{Imprécision du tracé.} À la règle et au crayon, une erreur d'un demi-millimètre sur le dessin fausse la lecture des coordonnées. La solution $\left(\dfrac{13}{7}\,;\dfrac{11}{7}\right)$ de l'exercice résolu est \emph{invisible} sur un graphique ordinaire.
\item \textbf{Coefficients non entiers.} Avec un système comme $\begin{cases} 0,37x - 1,2y = 4,5 \\ 2,9x + 0,83y = -1,7 \end{cases}$, le tracé précis des droites est quasi impossible à la main.
\item \textbf{Absence de preuve.} Une lecture graphique, même soignée, ne constitue \textbf{pas une démonstration} : elle fournit une conjecture que seul le calcul (substitution dans les deux équations, ou méthodes algébriques) peut confirmer.
\item \textbf{Cas ambigus.} Deux droites presque parallèles semblent se couper « très loin » ; deux droites presque confondues semblent n'en faire qu'une. Le dessin ne permet pas de trancher de manière fiable.
\end{itemize}
C'est pourquoi les sections suivantes développeront des \textbf{méthodes algébriques} (substitution, combinaison) qui donnent toujours la solution exacte.
\end{attention}

\begin{savaistu}[La géométrie analytique : quand l'algèbre rencontre la géométrie]
L'idée qu'une équation correspond à une courbe et que résoudre un système revient à chercher des intersections est l'une des plus grandes inventions des mathématiques : la \emph{géométrie analytique}, développée au XVII\textsuperscript{e} siècle par le Français René Descartes (c'est lui qui a donné son nom au « repère cartésien ») et Pierre de Fermat. Grâce à elle, un problème de géométrie (intersection de droites) devient un problème de calcul (résolution d'un système), et réciproquement. Aujourd'hui, ce principe est partout : le GPS qui te localise croise des informations de plusieurs satellites — c'est une résolution de systèmes ! — et les applications de cartographie utilisées pour naviguer dans Yaoundé ou planifier l'électrification rurale reposent sur ces mêmes idées.
\end{savaistu}

\begin{exercice}[À toi de jouer]
On considère le système $\begin{cases} 2x + y = 5 \\ x - y = 1 \end{cases}$.
\begin{enumerate}
\item Déterminer deux points de chaque droite, puis tracer les deux droites dans un repère (unité : 1 cm).
\item Lire graphiquement les coordonnées du point d'intersection.
\item Vérifier par le calcul que le couple lu est bien solution du système.
\item Sans nouveau tracé, dire combien de solutions possède le système $\begin{cases} 2x + y = 5 \\ 2x + y = 7 \end{cases}$. Justifier.
\end{enumerate}
\end{exercice}

\begin{corrige}[Exercice]
\begin{enumerate}
\item Droite $(d_1)$ : $2x + y = 5$. Si $x = 0$, $y = 5$ ; si $y = 0$, $x = 2,5$. Points $A(0\,;5)$ et $B(2,5\,;0)$.
Droite $(d_2)$ : $x - y = 1$. Si $x = 0$, $y = -1$ ; si $y = 0$, $x = 1$. Points $C(0\,;-1)$ et $D(1\,;0)$.
On trace les deux droites dans le repère.
\item Les droites se coupent au point $(2\,;1)$.
\item Vérification : $2\times 2 + 1 = 5$ \checkmark{} et $2 - 1 = 1$ \checkmark. Le couple $(2\,;1)$ est bien l'unique solution du système.
\item Les équations $2x + y = 5$ et $2x + y = 7$ ont les mêmes coefficients en $x$ et $y$ mais des seconds membres différents : les droites sont strictement parallèles. Le système n'a \textbf{aucune} solution (un même couple ne peut pas donner à $2x+y$ deux valeurs différentes).
\end{enumerate}
\end{corrige}

En résumé, l'approche graphique transforme un système en un problème d'intersection de droites : elle donne une vision claire des trois cas possibles (une solution, aucune, une infinité) et permet des lectures approchées. Mais son imprécision impose de la compléter par des méthodes de calcul exact : ce sera l'objet de la section suivante, consacrée à la substitution et à l'élimination.

\coursec{Méthodes algébriques : substitution et élimination (combinaison linéaire)}

Dans la section précédente, nous avons résolu des systèmes de deux équations du premier degré à deux inconnues par lecture \cle{graphique} : on trace les deux droites et on lit les coordonnées du point d'intersection. Cette méthode est parlante, mais elle a ses limites : si la solution est $(\dfrac{7}{3}\,;\,-\dfrac{5}{4})$, aucune lecture sur un graphique ne la donnera exactement ! Il nous faut donc des méthodes \cle{algébriques}, c'est-à-dire des calculs exacts qui fonctionnent à coup sûr, sans règle ni crayon. Il en existe deux grandes : la \cle{substitution} et l'\cle{élimination} (aussi appelée \emph{combinaison linéaire}). C'est ce que nous allons maîtriser maintenant, avec une exigence : \textbf{chaque ligne de calcul doit être justifiée}.

\courssub{La méthode de substitution}

\textbf{Idée intuitive.}\quad Imaginons qu'au marché de Mokolo, tu saches qu'un kilogramme de tomates coûte $100$ FCFA de plus qu'un kilogramme d'oignons. Si on te donne le prix des oignons, tu connais aussitôt celui des tomates : il suffit de \emph{remplacer} l'inconnu « prix des tomates » par « prix des oignons $+100$ ». La substitution, c'est exactement cela : l'une des équations nous donne une inconnue \emph{en fonction de} l'autre ; on \emph{substitue} (on remplace) cette expression dans l'autre équation, qui ne contient alors plus qu'\textbf{une seule inconnue} — et une équation à une inconnue, nous savons la résoudre depuis la classe de 4ème !

\begin{methode}[Résoudre un système par substitution]
Pour résoudre le système $\begin{cases} ax+by=c \\ a'x+b'y=c' \end{cases}$ par substitution :
\begin{enumerate}
\item \textbf{Isoler} : dans l'équation la plus simple (celle où une inconnue a pour coefficient $1$ ou $-1$), exprimer une inconnue en fonction de l'autre (par exemple $x = \ldots$).
\item \textbf{Substituer} : remplacer cette inconnue par son expression dans \emph{l'autre} équation (jamais dans celle dont elle provient !).
\item \textbf{Résoudre} : on obtient une équation du premier degré à une seule inconnue ; la résoudre.
\item \textbf{Remonter} : reporter la valeur trouvée dans l'expression de l'étape 1 pour obtenir la seconde inconnue.
\item \textbf{Vérifier} : remplacer le couple trouvé dans les \emph{deux} équations initiales, puis conclure par $S = \{(x\,;\,y)\}$.
\end{enumerate}
\end{methode}

\begin{exoresolu}[Résolution par substitution]
Résoudre dans $\mathbb{R}\times\mathbb{R}$ le système : $\begin{cases} 2x+3y=12 \\ x-y=1 \end{cases}$
\tcblower
\textbf{Étape 1 — Isoler.} Dans la deuxième équation, le coefficient de $x$ est $1$ : c'est le candidat idéal.
\[ x - y = 1 \iff x = y + 1. \]
\textbf{Étape 2 — Substituer.} On remplace $x$ par $y+1$ dans la \emph{première} équation :
\[ 2(y+1) + 3y = 12. \]
\textbf{Étape 3 — Résoudre.} On développe et on réduit :
\begin{align*}
2y + 2 + 3y &= 12 \\
5y &= 12 - 2 \\
5y &= 10 \\
y &= 2.
\end{align*}
\textbf{Étape 4 — Remonter.} $x = y + 1 = 2 + 1 = 3$.
\textbf{Étape 5 — Vérifier.} Dans la première équation : $2\times 3 + 3\times 2 = 6 + 6 = 12$ \checkmark. Dans la deuxième : $3 - 2 = 1$ \checkmark.
\textbf{Conclusion :} $S = \{(3\,;\,2)\}$.
\end{exoresolu}

La figure suivante résume l'enchaînement des étapes de cette résolution.

\begin{popfigure}
\begin{center}
\begin{tikzpicture}[scale=0.9]
\node[draw=popblue, rounded corners, fill=popblueL, align=left, inner sep=8pt] (s1) {\textbf{1. Isoler}\quad $x - y = 1 \Rightarrow x = y+1$};
\node[draw=poporange, rounded corners, fill=poporangeL, align=left, inner sep=8pt, below=0.35cm of s1] (s2) {\textbf{2. Substituer}\quad $2(y+1)+3y=12$};
\node[draw=popgreen, rounded corners, fill=popgreenL, align=left, inner sep=8pt, below=0.35cm of s2] (s3) {\textbf{3. Résoudre}\quad $5y=10 \Rightarrow y=2$};
\node[draw=poppurple, rounded corners, fill=poppurpleL, align=left, inner sep=8pt, below=0.35cm of s3] (s4) {\textbf{4. Remonter}\quad $x=2+1=3$ \quad \textbf{5. Vérifier}\quad $S=\{(3\,;\,2)\}$};
\draw[-{Stealth}, thick, popmuted] (s1) -- (s2);
\draw[-{Stealth}, thick, popmuted] (s2) -- (s3);
\draw[-{Stealth}, thick, popmuted] (s3) -- (s4);
\end{tikzpicture}
\end{center}
\legende{Les étapes de la méthode de substitution sur le système $\begin{cases}2x+3y=12\\x-y=1\end{cases}$.}
\end{popfigure}

\begin{attention}[Pièges de la substitution]
\begin{itemize}
\item \textbf{Ne jamais substituer dans l'équation de départ.} Si tu remplaces $x$ par $y+1$ dans $x - y = 1$, tu obtiens $(y+1)-y=1$, soit $1=1$ : une égalité vraie mais inutile, qui ne donne aucune valeur !
\item \textbf{Penser aux parenthèses.} Dans $2(y+1)$, le facteur $2$ multiplie \emph{toute} l'expression. Écrire $2y+1$ au lieu de $2y+2$ est l'erreur la plus fréquente.
\item \textbf{Ne pas oublier la seconde inconnue.} Trouver $y=2$ ne suffit pas : une solution d'un système à deux inconnues est un \emph{couple} $(x\,;\,y)$.
\end{itemize}
\end{attention}

\courssub{La méthode d'élimination (combinaison linéaire)}

\textbf{Idée intuitive.}\quad Deux amis pèsent ensemble $130$ kg et leur différence de poids est de $10$ kg. Si l'on « ajoute » les deux informations, la différence et la somme se combinent pour faire disparaître l'un des deux poids. L'élimination repose sur une idée toute simple : si deux quantités sont égales ($A = B$) et deux autres aussi ($C = D$), alors leurs sommes sont égales ($A + C = B + D$). En ajoutant membre à membre les deux équations du système — éventuellement après les avoir multipliées par des nombres bien choisis — on fait \emph{disparaître} (éliminer) une inconnue.

\begin{propriete}[Opérations licites sur un système]
On ne change pas les solutions d'un système lorsque :
\begin{itemize}
\item on multiplie (ou divise) \textbf{les deux membres} d'une équation par un même nombre non nul ;
\item on remplace une équation par la \textbf{somme membre à membre} des deux équations.
\end{itemize}
Ces deux opérations sont le fondement de la méthode d'élimination.
\end{propriete}

\begin{methode}[Résoudre un système par élimination]
\begin{enumerate}
\item \textbf{Choisir} l'inconnue à éliminer (celle dont les coefficients sont les plus faciles à rendre opposés).
\item \textbf{Multiplier} une ou les deux équations par des nombres bien choisis pour obtenir des coefficients \emph{opposés} (par exemple $+2y$ et $-2y$) devant cette inconnue.
\item \textbf{Additionner} membre à membre les deux équations : l'inconnue choisie disparaît.
\item \textbf{Résoudre} l'équation à une inconnue obtenue.
\item \textbf{Remonter} : reporter la valeur trouvée dans l'une des équations initiales pour calculer l'autre inconnue.
\item \textbf{Vérifier} dans les deux équations initiales, puis conclure.
\end{enumerate}
\end{methode}

\begin{exemplebox}[Élimination directe : les coefficients sont déjà opposés]
Résolvons $\begin{cases} 3x+2y=7 \quad (1)\\ 5x-2y=1 \quad (2) \end{cases}$

Les coefficients de $y$ sont $+2$ et $-2$ : déjà opposés ! Nul besoin de multiplier. On additionne membre à membre $(1) + (2)$ :
\[ (3x + 2y) + (5x - 2y) = 7 + 1 \iff 8x = 8 \iff x = 1. \]
On reporte $x=1$ dans l'équation $(1)$ : $3\times 1 + 2y = 7 \iff 2y = 4 \iff y = 2$.
\textbf{Vérification :} $(2)$ : $5\times 1 - 2\times 2 = 5 - 4 = 1$ \checkmark.
\textbf{Conclusion :} $S = \{(1\,;\,2)\}$.
\end{exemplebox}

\begin{exoresolu}[Élimination avec multiplication préalable]
Résoudre dans $\mathbb{R}\times\mathbb{R}$ : $\begin{cases} 2x+3y=8 \quad (1)\\ 3x-y=1 \quad (2) \end{cases}$
\tcblower
Aucun coefficient n'est opposé à un autre. Choisissons d'éliminer $y$ : il suffit de multiplier l'équation $(2)$ par $3$ pour obtenir $-3y$, opposé du $+3y$ de $(1)$.
\[ (2) \times 3 : \quad 9x - 3y = 3 \quad (2') \]
On additionne $(1) + (2')$ :
\begin{align*}
(2x + 3y) + (9x - 3y) &= 8 + 3 \\
11x &= 11 \\
x &= 1.
\end{align*}
On reporte $x = 1$ dans $(2)$ : $3\times 1 - y = 1 \iff y = 2$.
\textbf{Vérification :} $(1)$ : $2\times 1 + 3\times 2 = 2 + 6 = 8$ \checkmark ; $(2)$ : $3 - 2 = 1$ \checkmark.
\textbf{Conclusion :} $S = \{(1\,;\,2)\}$.
\end{exoresolu}

\begin{attention}[Multiplier les DEUX membres !]
Quand on multiplie l'équation $3x - y = 1$ par $3$, on obtient $9x - 3y = \mathbf{3}$. L'erreur classique consiste à multiplier le membre de gauche et à oublier le membre de droite (écrire $9x - 3y = 1$). Rappelle-toi : une équation est une balance ; ce que tu fais à gauche, tu le fais à droite, sinon tu fausses tout le calcul.
\end{attention}

\courssub{Quelle méthode choisir ?}

Les deux méthodes donnent \emph{toujours} le même résultat ; seule l'efficacité change.

\begin{aretenir}[Choisir sa méthode]
\begin{itemize}
\item Si une inconnue a pour coefficient $1$ ou $-1$ dans l'une des équations (par exemple $x - y = 1$ ou $y = 2x - 5$), la \textbf{substitution} est la plus rapide : on isole sans fraction.
\item Si tous les coefficients sont différents de $\pm 1$ et que des coefficients sont déjà opposés ou faciles à rendre opposés (comme $+2y$ et $-2y$), l'\textbf{élimination} est préférable : elle évite les fractions.
\item Dans le doute, l'élimination est souvent la plus sûre, car la substitution avec des coefficients quelconques fait apparaître des fractions sources d'erreurs.
\end{itemize}
\end{aretenir}

\courssub{Cas particuliers : système impossible et système indéterminé}

Que se passe-t-il lorsque les deux droites du système sont parallèles (aucune solution) ou confondues (une infinité de solutions) ? Le calcul algébrique le détecte tout seul : au moment de l'élimination, \emph{les deux inconnues disparaissent en même temps}.

\begin{exemplebox}[Système impossible : on aboutit à $0 = k$ avec $k \neq 0$]
Résolvons $\begin{cases} x + 2y = 5 \quad (1)\\ 2x + 4y = 7 \quad (2) \end{cases}$

Multiplions $(1)$ par $2$ : $2x + 4y = 10 \quad (1')$. Soustrayons $(2)$ de $(1')$ :
\[ (2x + 4y) - (2x + 4y) = 10 - 7 \iff 0 = 3. \]
Cette égalité est \textbf{fausse}, quelles que soient les valeurs de $x$ et $y$ : aucun couple ne peut vérifier le système. On conclut : $S = \emptyset$ (système \cle{impossible}). Graphiquement, les deux droites sont strictement parallèles.
\end{exemplebox}

\begin{exemplebox}[Système indéterminé : on aboutit à $0 = 0$]
Résolvons $\begin{cases} x + 2y = 5 \quad (1)\\ 2x + 4y = 10 \quad (2) \end{cases}$

Multiplions $(1)$ par $2$ : $2x + 4y = 10$ : c'est exactement l'équation $(2)$ ! En soustrayant, on obtient $0 = 0$, égalité \textbf{toujours vraie}. Les deux équations disent la même chose : tout couple vérifiant $x + 2y = 5$ est solution. Il y a une \textbf{infinité de solutions} (système \cle{indéterminé}) : $S = \{(5 - 2y\,;\,y),\ y \in \mathbb{R}\}$. Graphiquement, les deux droites sont confondues.
\end{exemplebox}

\begin{aretenir}[Les trois issues possibles d'un système]
Après avoir tenté d'éliminer une inconnue (par combinaison linéaire), trois cas peuvent se présenter :
\begin{itemize}
\item \textbf{Solution unique :} on obtient une équation du type $ax = b$ avec $a \neq 0$ (une inconnue reste). On calcule $x$, puis on remonte pour trouver $y$.
\item \textbf{Aucune solution (système impossible) :} les deux inconnues disparaissent et on obtient une égalité fausse $0 = k$ avec $k \neq 0$. Exemple : $0 = 3$. Les droites sont parallèles strictes : $S = \emptyset$.
\item \textbf{Infinité de solutions (système indéterminé) :} les deux inconnues disparaissent et on obtient une égalité vraie $0 = 0$. Les deux équations sont proportionnelles (même droite). On exprime l'ensemble des solutions en fonction d'un paramètre : $S = \{(x\,;\,y) \in \mathbb{R}^2 \mid ax+by=c\}$.
\end{itemize}
\end{aretenir}

\courssub{Rédaction formelle : la présentation attendue au BEPC}

Au BEPC, la présentation compte autant que le résultat. Voici le modèle à reproduire.

\begin{methode}[Rédaction type d'une résolution de système]
\begin{enumerate}
\item Numéroter les équations $(1)$ et $(2)$.
\item Annoncer la méthode choisie : « Résolvons par élimination » ou « par substitution ».
\item Indiquer chaque opération effectuée : « $(2)\times 3$ », « $(1)+(2')$ », « on remplace $x$ par $y+1$ dans $(1)$ ».
\item Justifier chaque ligne de calcul par une équivalence ($\iff$) ou une implication claire.
\item Vérifier explicitement le couple trouvé dans les deux équations initiales.
\item Conclure par une phrase : « Le système admet pour unique solution le couple $(1\,;\,2)$ : $S = \{(1\,;\,2)\}$. »
\end{enumerate}
\end{methode}

\begin{exercice}[À toi de jouer]
\begin{enumerate}
\item Résoudre par substitution : $\begin{cases} y = 2x - 1 \\ 3x + y = 9 \end{cases}$
\item Résoudre par élimination : $\begin{cases} 4x + 3y = 10 \\ 2x - 3y = 2 \end{cases}$
\item Résoudre et interpréter : $\begin{cases} 2x - y = 3 \\ 4x - 2y = 5 \end{cases}$
\end{enumerate}
\end{exercice}

\begin{corrige}[Exercice]
\textbf{1.} On remplace $y$ par $2x-1$ dans la deuxième équation : $3x + (2x - 1) = 9 \iff 5x = 10 \iff x = 2$. Alors $y = 2\times 2 - 1 = 3$. Vérification : $3\times 2 + 3 = 9$ \checkmark. $S = \{(2\,;\,3)\}$.

\textbf{2.} Les coefficients de $y$ sont opposés ($+3$ et $-3$). On additionne : $6x = 12 \iff x = 2$. Dans la deuxième équation : $2\times 2 - 3y = 2 \iff 3y = 2 \iff y = \dfrac{2}{3}$. Vérification : $4\times 2 + 3\times\dfrac{2}{3} = 8 + 2 = 10$ \checkmark. $S = \left\{\left(2\,;\,\dfrac{2}{3}\right)\right\}$.

\textbf{3.} Multiplions la première équation par $2$ : $4x - 2y = 6$. En soustrayant la deuxième : $0 = 1$, égalité fausse. Le système est impossible : $S = \emptyset$ (droites strictement parallèles).
\end{corrige}

\begin{savaistu}[De l'élimination à la méthode de Gauss]
La méthode d'élimination que tu viens d'apprendre est l'ancêtre direct de la célèbre \textbf{méthode du pivot de Gauss}, du nom du mathématicien allemand Carl Friedrich Gauss (1777--1855). Elle permet de résoudre des systèmes de 3, 10 ou même des millions d'équations ! C'est elle qui tourne au cœur des ordinateurs : prévisions météorologiques, calculs de structures de ponts, moteurs de recherche, traitement du signal dans ton téléphone MTN ou Orange... Chaque fois qu'un réseau mobile optimise ses antennes à Yaoundé ou à Douala, des systèmes géants sont résolus par des variantes de ta méthode d'élimination. Un outil de 3ème, mais une puissance mondiale !
\end{savaistu}

\begin{aretenir}[L'essentiel de la section]
\begin{itemize}
\item \textbf{Substitution} : isoler une inconnue (coefficient $1$ ou $-1$), la remplacer dans l'autre équation, résoudre, remonter, vérifier.
\item \textbf{Élimination} : multiplier pour créer des coefficients opposés, additionner membre à membre, résoudre, remonter, vérifier.
\item Multiplier une équation, c'est multiplier \textbf{les deux membres}.
\item $0 = k$ ($k\neq 0$) : aucune solution ; $0 = 0$ : infinité de solutions.
\item Toujours \textbf{vérifier} le couple trouvé dans les deux équations initiales et conclure par $S = \{(x\,;\,y)\}$.
\end{itemize}
\end{aretenir}

\coursec{Discussion du nombre de solutions : déterminant et cas particuliers}

Dans la section précédente, tu as appris à résoudre un système de deux équations à deux inconnues par \cle{substitution} et par \cle{combinaison linéaire} (élimination). Tu as peut-être remarqué un phénomène curieux : certaines résolutions aboutissent à une absurdité comme $0 = 5$, d'autres à une évidence comme $0 = 0$, et la plupart donnent un couple solution unique. Faut-il vraiment mener tous les calculs pour savoir dans quel cas on se trouve ? \textbf{Non !} Il existe un outil remarquable, un simple nombre calculé à partir des coefficients du système, qui permet de \emph{prédire} à l'avance le nombre de solutions : le \cle{déterminant}. C'est l'objet de cette section.

\courssub{Écriture générale d'un système et déterminant}

Considérons un système de deux équations du premier degré à deux inconnues $x$ et $y$, écrit sous sa forme générale :
\[
\begin{cases} ax + by = e \\ cx + dy = f \end{cases}
\]
où $a$, $b$, $c$, $d$, $e$, $f$ sont des nombres réels connus (avec $a$ et $b$ non tous les deux nuls, de même pour $c$ et $d$).

On peut regrouper les coefficients des inconnues dans un tableau à deux lignes et deux colonnes, appelé \cle{matrice} :
\[
A = \begin{pmatrix} a & b \\ c & d \end{pmatrix}
\]
La première ligne contient les coefficients de la première équation, la deuxième ligne ceux de la deuxième équation ; la première colonne contient les coefficients de $x$, la deuxième ceux de $y$. Le système s'écrit alors symboliquement $A \cdot X = B$, avec $X = \begin{pmatrix} x \\ y \end{pmatrix}$ le couple inconnu et $B = \begin{pmatrix} e \\ f \end{pmatrix}$ le second membre. Cette écriture compacte est très utilisée en sciences ; en 3ème, nous l'utilisons surtout pour introduire le nombre qui va tout décider.

\begin{definition}[Déterminant d'un système]
Le \cle{déterminant} du système $\begin{cases} ax + by = e \\ cx + dy = f \end{cases}$ est le nombre réel noté $\Delta$ (lettre grecque « delta ») défini par :
\[
\Delta = ad - bc
\]
C'est le produit des coefficients de la \emph{diagonale principale} ($a$ et $d$) moins le produit des coefficients de l'\emph{autre diagonale} ($b$ et $c$).
\end{definition}

\begin{exemplebox}[Premiers calculs de déterminants]
\begin{itemize}
\item Pour $\begin{cases} 3x + 2y = 7 \\ x - 5y = 4 \end{cases}$ : $\Delta = 3 \times (-5) - 2 \times 1 = -15 - 2 = -17$.
\item Pour $\begin{cases} 2x + 4y = 6 \\ x + 2y = 3 \end{cases}$ : $\Delta = 2 \times 2 - 4 \times 1 = 4 - 4 = 0$.
\item Pour $\begin{cases} x + 2y = 4 \\ 2x + 4y = 6 \end{cases}$ : $\Delta = 1 \times 4 - 2 \times 2 = 4 - 4 = 0$.
\end{itemize}
Observe bien : dans les deux derniers cas, les coefficients de $x$ et $y$ sont \emph{proportionnels} d'une équation à l'autre, et le déterminant est nul. Ce n'est pas un hasard, comme nous allons le voir.
\end{exemplebox}

\courssub{Les déterminants $\Delta_x$ et $\Delta_y$}

Pour énoncer les formules de résolution, on introduit deux autres déterminants, obtenus en remplaçant une colonne de la matrice par le second membre.

\begin{definition}[Déterminants associés]
On pose :
\[
\Delta_x = \begin{vmatrix} e & b \\ f & d \end{vmatrix} = ed - bf
\qquad \text{et} \qquad
\Delta_y = \begin{vmatrix} a & e \\ c & f \end{vmatrix} = af - ec
\]
$\Delta_x$ s'obtient en remplaçant la colonne des coefficients de $x$ par le second membre $\begin{pmatrix} e \\ f \end{pmatrix}$ ; $\Delta_y$ s'obtient en remplaçant la colonne des coefficients de $y$.
\end{definition}

\begin{experience}[Démonstration guidée : d'où viennent ces formules ?]
Partons du système $\begin{cases} ax + by = e \quad (1) \\ cx + dy = f \quad (2) \end{cases}$ et utilisons la méthode de combinaison linéaire vue à la section 3.

\textbf{Étape 1 : éliminer $y$.} Multiplions l'équation (1) par $d$ et l'équation (2) par $b$ :
\[
adx + bdy = ed \qquad \text{et} \qquad bcx + bdy = bf
\]
En soustrayant membre à membre, les termes en $y$ disparaissent :
\[
(ad - bc)\,x = ed - bf \quad \text{c'est-à-dire} \quad \Delta \cdot x = \Delta_x
\]

\textbf{Étape 2 : éliminer $x$.} Multiplions (1) par $c$ et (2) par $a$ :
\[
acx + bcy = ec \qquad \text{et} \qquad acx + ady = af
\]
En soustrayant : $(ad - bc)\,y = af - ec$, c'est-à-dire $\Delta \cdot y = \Delta_y$.

\textbf{Étape 3 : conclusion.} Tout couple $(x\,;y)$ solution du système vérifie donc obligatoirement :
\[
\Delta \cdot x = \Delta_x \qquad \text{et} \qquad \Delta \cdot y = \Delta_y
\]
Toute la discussion va se jouer sur ces deux égalités : peut-on diviser par $\Delta$ ? Cela dépend de sa nullité !
\end{experience}

\courssub{Premier cas : $\Delta \neq 0$, le système de Cramer}

Si $\Delta \neq 0$, on peut diviser les deux égalités $\Delta \cdot x = \Delta_x$ et $\Delta \cdot y = \Delta_y$ par $\Delta$. On obtient immédiatement le théorème fondamental de cette leçon.

\begin{propriete}[Théorème de Cramer]
Si le déterminant $\Delta = ad - bc$ est \textbf{non nul}, alors le système $\begin{cases} ax + by = e \\ cx + dy = f \end{cases}$ admet une \textbf{unique solution}, donnée par les \emph{formules de Cramer} :
\[
x = \dfrac{\Delta_x}{\Delta} = \dfrac{ed - bf}{ad - bc}
\qquad \text{et} \qquad
y = \dfrac{\Delta_y}{\Delta} = \dfrac{af - ec}{ad - bc}
\]
Géométriquement, les deux droites représentant les équations sont \cle{sécantes} : elles se coupent en un seul point, le couple solution.
\end{propriete}

\begin{savaistu}[Gabriel Cramer]
Gabriel Cramer (1704--1752) était un mathématicien suisse. Ses célèbres formules, publiées en 1750, permettent de résoudre les systèmes d'équations à l'aide de déterminants. Aujourd'hui encore, les ordinateurs qui pilotent les réseaux de téléphonie ou les prévisions météo utilisent des généralisations de ces idées pour résoudre des systèmes comportant des millions d'inconnues !
\end{savaistu}

\begin{exoresolu}[Résolution par les formules de Cramer]
Résoudre dans $\mathbb{R} \times \mathbb{R}$ le système : $\begin{cases} 3x + 2y = 7 \\ x - 5y = 4 \end{cases}$
\tcblower
\textbf{Solution détaillée.}

\textbf{Étape 1 : calcul de $\Delta$.}
\[
\Delta = 3 \times (-5) - 2 \times 1 = -15 - 2 = -17
\]
$\Delta = -17 \neq 0$ : le système admet une unique solution (système de Cramer).

\textbf{Étape 2 : calcul de $\Delta_x$ et $\Delta_y$.}
\[
\Delta_x = 7 \times (-5) - 2 \times 4 = -35 - 8 = -43
\]
\[
\Delta_y = 3 \times 4 - 7 \times 1 = 12 - 7 = 5
\]

\textbf{Étape 3 : application des formules.}
\[
x = \dfrac{\Delta_x}{\Delta} = \dfrac{-43}{-17} = \dfrac{43}{17}
\qquad ; \qquad
y = \dfrac{\Delta_y}{\Delta} = \dfrac{5}{-17} = -\dfrac{5}{17}
\]

\textbf{Vérification} dans la première équation : $3 \times \dfrac{43}{17} + 2 \times \left(-\dfrac{5}{17}\right) = \dfrac{129 - 10}{17} = \dfrac{119}{17} = 7$. \checkmark

\textbf{Conclusion :} $S = \left\lbrace \left(\dfrac{43}{17}\,;\, -\dfrac{5}{17}\right) \right\rbrace$.
\end{exoresolu}

\courssub{Deuxième cas : $\Delta = 0$, les cas particuliers}

Revenons aux égalités fondamentales obtenues dans la démonstration guidée :
\[
\Delta \cdot x = \Delta_x \qquad \text{et} \qquad \Delta \cdot y = \Delta_y
\]
Si $\Delta = 0$, ces égalités deviennent $0 \times x = \Delta_x$ et $0 \times y = \Delta_y$. Or $0$ fois n'importe quel nombre vaut toujours $0$. Deux situations sont alors possibles.

\textbf{Sous-cas 1 : $\Delta = 0$ et ($\Delta_x \neq 0$ ou $\Delta_y \neq 0$).}

L'égalité $0 \times x = \Delta_x$ est alors \emph{impossible} (par exemple $0 \times x = 5$ n'a aucune solution). Le système n'admet \textbf{aucune solution}. Géométriquement, les deux droites sont \cle{parallèles et distinctes} : elles ne se rencontrent jamais.

\textbf{Sous-cas 2 : $\Delta = 0$ et $\Delta_x = \Delta_y = 0$.}

Les égalités $0 \times x = 0$ et $0 \times y = 0$ sont toujours vraies : elles n'imposent aucune condition. Les deux équations sont en fait \emph{proportionnelles} (l'une est un multiple de l'autre) : elles représentent la \textbf{même droite}. Le système admet une \textbf{infinité de solutions} : tous les couples $(x\,;y)$ vérifiant l'une des deux équations. Les droites sont \cle{confondues}.

\begin{propriete}[Bilan de la discussion]
Pour le système $\begin{cases} ax + by = e \\ cx + dy = f \end{cases}$, on calcule $\Delta = ad - bc$, $\Delta_x = ed - bf$ et $\Delta_y = af - ec$ :
\begin{center}
\begin{tabularx}{\linewidth}{|l|X|X|}
\hline
\rowcolor{popblueL} \textbf{Condition} & \textbf{Nombre de solutions} & \textbf{Position des droites} \\
\hline
$\Delta \neq 0$ & Une solution unique & Droites sécantes \\
\hline
$\Delta = 0$ et ($\Delta_x \neq 0$ ou $\Delta_y \neq 0$) & Aucune solution & Droites parallèles distinctes \\
\hline
$\Delta = 0$ et $\Delta_x = \Delta_y = 0$ & Infinité de solutions & Droites confondues \\
\hline
\end{tabularx}
\end{center}
\end{propriete}

\begin{propriete}[Lien avec le parallélisme]
$\Delta = ad - bc = 0$ si et seulement si les coefficients des inconnues sont proportionnels ($\dfrac{a}{c} = \dfrac{b}{d}$ lorsque $c$ et $d$ sont non nuls), c'est-à-dire si et seulement si les deux droites du système sont \textbf{parallèles ou confondues}. En effet, les vecteurs directeurs $\vec{u}(-b\,;a)$ et $\vec{v}(-d\,;c)$ des deux droites sont colinéaires exactement lorsque $(-b)c - a(-d) = ad - bc = 0$.
\end{propriete}

\begin{exemplebox}[Exemple chiffré : droites confondues]
Considérons le système $\begin{cases} 2x + 4y = 6 \\ x + 2y = 3 \end{cases}$.

Calculons les trois déterminants :
\[
\Delta = 2 \times 2 - 4 \times 1 = 4 - 4 = 0
\]
\[
\Delta_x = 6 \times 2 - 4 \times 3 = 12 - 12 = 0
\qquad ; \qquad
\Delta_y = 2 \times 3 - 6 \times 1 = 6 - 6 = 0
\]
Comme $\Delta = \Delta_x = \Delta_y = 0$, le système admet une \textbf{infinité de solutions}. C'est logique : en divisant la première équation par $2$, on obtient exactement la deuxième ! Les deux équations représentent la même droite. L'ensemble des solutions est :
\[
S = \left\lbrace (x\,;y) \in \mathbb{R}^2 \; / \; x + 2y = 3 \right\rbrace = \left\lbrace (3 - 2y\,;\,y),\; y \in \mathbb{R} \right\rbrace
\]
Par exemple, $(3\,;0)$, $(1\,;1)$, $(-1\,;2)$ sont toutes des solutions.
\end{exemplebox}

\begin{popfigure}
\begin{tikzpicture}[scale=0.7]
\draw[->] (-1,0)--(5,0) node[right]{$x$};
\draw[->] (0,-1)--(0,4) node[above]{$y$};
\draw[thick,red,dashed] (0,2)--(4,0) node[midway,above right]{$x+2y=4$};
\draw[thick,blue,dashed] (0,1.5)--(3,0) node[midway,above right]{$2x+4y=6$};
\node at (2,3) {Parallèles distinctes : $\Delta=0$, $\Delta_x \neq 0$};
\end{tikzpicture}
\legende{Le système $\begin{cases} x+2y=4 \\ 2x+4y=6 \end{cases}$ : $\Delta = 1 \times 4 - 2 \times 2 = 0$ et $\Delta_x = 4 \times 4 - 2 \times 6 = 4 \neq 0$. Les deux droites sont parallèles et distinctes : aucune solution.}
\end{popfigure}

\courssub{Méthode générale de discussion}

\begin{methode}[Discuter le nombre de solutions d'un système]
\begin{enumerate}
\item Écrire le système sous la forme $\begin{cases} ax + by = e \\ cx + dy = f \end{cases}$ (réduire et ordonner si nécessaire).
\item Calculer le déterminant $\Delta = ad - bc$.
\item \textbf{Si $\Delta \neq 0$} : conclure « solution unique » et la calculer par $x = \dfrac{\Delta_x}{\Delta}$, $y = \dfrac{\Delta_y}{\Delta}$.
\item \textbf{Si $\Delta = 0$} : calculer $\Delta_x = ed - bf$ et $\Delta_y = af - ec$.
\begin{itemize}
\item Si l'un des deux est non nul : conclure « aucune solution » (droites parallèles distinctes).
\item Si les deux sont nuls : conclure « infinité de solutions » (droites confondues) et décrire l'ensemble $S$.
\end{itemize}
\end{enumerate}
\end{methode}

\begin{attention}[Pièges fréquents]
\begin{itemize}
\item \textbf{Ne jamais diviser par $\Delta$ s'il est nul !} Les formules de Cramer $x = \dfrac{\Delta_x}{\Delta}$ et $y = \dfrac{\Delta_y}{\Delta}$ ne s'appliquent \emph{que} lorsque $\Delta \neq 0$. Diviser par zéro est interdit en mathématiques.
\item Attention à l'\textbf{ordre des termes} dans le déterminant : $\Delta = ad - bc$, et non $ac - bd$. C'est le produit des termes de la diagonale principale moins le produit des termes de l'autre diagonale.
\item Avant de calculer $\Delta$, il faut avoir \textbf{ordonné} le système : les termes en $x$ d'abord, puis les termes en $y$, puis le second membre. Pour $\begin{cases} 2y + x = 5 \\ 3x = y - 1 \end{cases}$, on réécrit d'abord $\begin{cases} x + 2y = 5 \\ 3x - y = -1 \end{cases}$.
\item « $\Delta = 0$ » ne signifie pas automatiquement « aucune solution » : il faut examiner $\Delta_x$ et $\Delta_y$ pour trancher entre « aucune » et « une infinité ».
\end{itemize}
\end{attention}

\begin{savaistu}[Le déterminant et l'aire du parallélogramme]
Le déterminant n'apparaît pas seulement dans les systèmes ! En géométrie, si deux vecteurs $\vec{u}(a\,;b)$ et $\vec{v}(c\,;d)$ engendrent un parallélogramme, alors l'\textbf{aire} de ce parallélogramme est égale à la valeur absolue du déterminant : $\mathcal{A} = |ad - bc|$. Ainsi, dire que $\Delta = 0$ revient à dire que le parallélogramme est « aplati » : les deux vecteurs sont colinéaires, ils pointent dans la même direction (ou des directions opposées). C'est exactement pour cela que $\Delta = 0$ correspond au parallélisme des droites du système. Les mathématiciens disent que le déterminant mesure « l'aire algébrique » : une même notion relie l'algèbre et la géométrie !
\end{savaistu}

\begin{exoresolu}[Discussion complète]
Sans résoudre entièrement, déterminer le nombre de solutions de chacun des systèmes suivants, puis résoudre celui qui admet une solution unique :
\[
(S_1) : \begin{cases} 5x - y = 3 \\ 2x + 3y = 8 \end{cases}
\qquad
(S_2) : \begin{cases} x + 2y = 4 \\ 2x + 4y = 6 \end{cases}
\]
\tcblower
\textbf{Solution détaillée.}

\textbf{Pour $(S_1)$ :} $\Delta = 5 \times 3 - (-1) \times 2 = 15 + 2 = 17 \neq 0$.

Le système admet une unique solution. Calculons :
\[
\Delta_x = 3 \times 3 - (-1) \times 8 = 9 + 8 = 17
\qquad ; \qquad
\Delta_y = 5 \times 8 - 3 \times 2 = 40 - 6 = 34
\]
\[
x = \dfrac{17}{17} = 1 \qquad ; \qquad y = \dfrac{34}{17} = 2
\]
Vérification : $5(1) - 2 = 3$ \checkmark{} et $2(1) + 3(2) = 8$ \checkmark. Donc $S_1 = \lbrace (1\,;2) \rbrace$.

\textbf{Pour $(S_2)$ :} $\Delta = 1 \times 4 - 2 \times 2 = 4 - 4 = 0$. On calcule alors :
\[
\Delta_x = 4 \times 4 - 2 \times 6 = 16 - 12 = 4 \neq 0
\]
Comme $\Delta = 0$ et $\Delta_x \neq 0$, le système $(S_2)$ n'admet \textbf{aucune solution} : les deux droites sont parallèles et distinctes (voir la figure ci-dessus). On écrit $S_2 = \varnothing$.
\end{exoresolu}

\begin{exercice}[À toi de jouer]
Pour chaque système, calculer $\Delta$, puis $\Delta_x$ et $\Delta_y$ si nécessaire, et conclure sur le nombre de solutions (sans chercher les solutions) :
\[
\text{a) } \begin{cases} 4x - 3y = 1 \\ 2x + y = 7 \end{cases}
\qquad
\text{b) } \begin{cases} 6x - 9y = 12 \\ 2x - 3y = 4 \end{cases}
\qquad
\text{c) } \begin{cases} 3x + 6y = 9 \\ x + 2y = 5 \end{cases}
\]
\end{exercice}

\begin{corrige}[Exercice : discussion]
\textbf{a)} $\Delta = 4 \times 1 - (-3) \times 2 = 4 + 6 = 10 \neq 0$ : solution unique.

\textbf{b)} $\Delta = 6 \times (-3) - (-9) \times 2 = -18 + 18 = 0$. Puis $\Delta_x = 12 \times (-3) - (-9) \times 4 = -36 + 36 = 0$ et $\Delta_y = 6 \times 4 - 12 \times 2 = 24 - 24 = 0$. Les trois déterminants sont nuls : infinité de solutions (droites confondues ; la première équation est le triple de la deuxième).

\textbf{c)} $\Delta = 3 \times 2 - 6 \times 1 = 6 - 6 = 0$. Puis $\Delta_x = 9 \times 2 - 6 \times 5 = 18 - 30 = -12 \neq 0$. Aucune solution : droites parallèles distinctes.
\end{corrige}

\begin{aretenir}[L'essentiel de la section]
\begin{itemize}
\item Le \cle{déterminant} du système $\begin{cases} ax + by = e \\ cx + dy = f \end{cases}$ est $\Delta = ad - bc$.
\item Si $\Delta \neq 0$ : \textbf{solution unique} donnée par les formules de Cramer $x = \dfrac{\Delta_x}{\Delta}$ et $y = \dfrac{\Delta_y}{\Delta}$, avec $\Delta_x = ed - bf$ et $\Delta_y = af - ec$ (droites sécantes).
\item Si $\Delta = 0$ et ($\Delta_x \neq 0$ ou $\Delta_y \neq 0$) : \textbf{aucune solution} (droites parallèles distinctes).
\item Si $\Delta = 0$ et $\Delta_x = \Delta_y = 0$ : \textbf{infinité de solutions} (droites confondues).
\item $\Delta = 0$ traduit géométriquement le \cle{parallélisme} (ou la confusion) des deux droites : leurs vecteurs directeurs sont colinéaires.
\end{itemize}
\end{aretenir}

Tu es maintenant capable, grâce à trois petits calculs de déterminants, de prédire le destin de n'importe quel système de deux équations à deux inconnues avant même de le résoudre. Dans la section suivante, nous mettrons tous ces outils au service de la résolution de problèmes concrets de la vie courante : tu verras que le déterminant t'aidera à vérifier rapidement si un problème est bien posé !

\coursec{Modélisation et résolution de problèmes concrets}

Après avoir maîtrisé les outils algébriques (substitution, élimination, règle de Cramer) et compris la classification des systèmes (déterminé, indéterminé, impossible) grâce au déterminant, nous voici au cœur de la puissance des mathématiques : \textbf{la modélisation}. Une équation ou un système n'est pas une fin en soi ; c'est le \emph{langage} qui permet de traduire une situation réelle pour en extraire la solution. Dans cette section, nous allons apprendre à passer du texte au système, le résoudre avec méthode, et surtout \textbf{redonner du sens} à la réponse trouvée.

\courssub{La démarche de modélisation : du concret à l'abstrait}

Modéliser, c'est construire un \emph{modèle mathématique} (ici, un système de deux équations du premier degré à deux inconnues) qui reflète fidèlement les contraintes d'un problème. Cette démarche suit toujours les mêmes étapes rigoureuses.

\begin{methode}[Les 5 étapes de la modélisation par un système linéaire]
\begin{enumerate}
    \item \textbf{Analyser et choisir les inconnues.} Lire attentivement le problème. Identifier les \emph{deux} quantités inconnues recherchées. Les nommer explicitement (ex : « Soit $x$ le nombre de beignets et $y$ le nombre de boissons »). Préciser leur unité et leur ensemble de définition naturel (souvent $\mathbb{N}$ ou $\mathbb{R}_+$).
    \item \textbf{Traduire les contraintes en équations.} Repérer \emph{deux} relations indépendantes entre ces quantités (prix total, quantité totale, périmètre, différence, rapport, etc.). Les écrire sous forme d'équations du premier degré : $ax + by = c$.
    \item \textbf{Former et résoudre le système.} Écrire le système $\begin{cases} a_1x+b_1y=c_1 \\ a_2x+b_2y=c_2 \end{cases}$. Choisir la méthode la plus adaptée (substitution si une inconnue est isolée facilement, élimination si les coefficients se prêtent à une annulation, Cramer si les calculs sont simples). Effectuer les calculs avec soin.
    \item \textbf{Interpréter la solution dans le contexte.} Vérifier que le couple $(x;y)$ obtenu a du sens \emph{réel} : est-il dans $\mathbb{N}$ (objets comptables) ? Est-il positif ? Respecte-t-il les contraintes implicites (ex : « au moins 3 articles ») ? Si le système est indéterminé ou impossible, que signifie cela pour le problème ?
    \item \textbf{Rédiger la réponse finale.} Écrire une phrase complète, en français correct, reprenant les mots du problème, avec les unités. Exemple : « Maman Ngo a acheté 6 beignets et 3 boissons. »
\end{enumerate}
\end{methode}

\begin{attention}[Pièges classiques de la modélisation]
\begin{itemize}
    \item \textbf{Confondre les inconnues :} Ne pas écrire « Soit $x$ le prix d'un beignet » si la question demande « Combien de beignets ? ».
    \item \textbf{Oublier l'ensemble de définition :} Trouver $x = 3,5$ pour un nombre d'objets est un signal d'alerte : erreur de calcul, de traduction, ou problème mal posé (données incohérentes).
    \item \textbf{Équations redondantes ou contradictoires :} Si les deux phrases du problème disent la même chose (système indéterminé) ou des choses impossibles ensemble (système impossible), le problème n'a pas de solution unique ou n'a pas de sens.
    \item \textbf{Mauvaise élimination :} Multiplier une équation par un nombre sans multiplier \emph{tous} les termes (y compris le second membre).
\end{itemize}
\end{attention}

\courssub{Exemple type complet : Le budget de Maman Ngo (Achats)}

C'est le modèle classique de la « double contrainte » : une contrainte \emph{quantitative} (nombre total d'articles) et une contrainte \emph{financière} (coût total).

\begin{exoresolu}[Problème : Beignets et Boissons]
\textbf{Énoncé.} Maman Ngo va au marché pour la rentrée. Elle achète des beignets à 150 FCFA l'unité et des boissons à 200 FCFA l'unité. Elle dépense au total 1500 FCFA et repart avec 9 articles au total. Combien a-t-elle acheté de beignets et de boissons ?

\tcblower

\textbf{Solution détaillée.}

\noindent \textbf{1. Choix des inconnues.}
Soit $x$ le nombre de beignets achetés.
Soit $y$ le nombre de boissons achetées.
$x$ et $y$ sont des entiers naturels ($x, y \in \mathbb{N}$).

\noindent \textbf{2. Traduction des contraintes.}
\begin{itemize}
    \item Contrainte quantitative (nombre d'articles) : $x + y = 9$.
    \item Contrainte financière (coût total) : $150x + 200y = 1500$.
\end{itemize}

\noindent \textbf{3. Formation et résolution du système.}
On obtient le système :
\[
(S) : \begin{cases}
x + y = 9 & (L_1)\\
150x + 200y = 1500 & (L_2)
\end{cases}
\]

\emph{Choix de la méthode :} La première équation est très simple ($x = 9 - y$). La \textbf{méthode par substitution} est idéale.

\begin{align*}
\text{De } (L_1) : \quad & x = 9 - y \\
\text{On substitue dans } (L_2) : \quad & 150(9 - y) + 200y = 1500 \\
& 1350 - 150y + 200y = 1500 \\
& 50y = 1500 - 1350 \\
& 50y = 150 \\
& y = \frac{150}{50} = 3
\end{align*}

On retrouve $x$ :
\[
x = 9 - y = 9 - 3 = 6
\]

La solution du système est le couple $(6; 3)$.

\noindent \textbf{4. Interprétation.}
$x = 6$ et $y = 3$ sont des entiers naturels positifs. Cela a un sens concret : on ne peut pas acheter un nombre négatif ou décimal d'articles. La solution est \textbf{acceptable}.

\noindent \textbf{5. Vérification (recommandée).}
\begin{itemize}
    \item Nombre total : $6 + 3 = 9$ ✓
    \item Coût total : $6 \times 150 + 3 \times 200 = 900 + 600 = 1500$ FCFA ✓
\end{itemize}

\noindent \textbf{6. Réponse finale.}
\fbox{\parbox{\linewidth}{\centering \textbf{Maman Ngo a acheté 6 beignets et 3 boissons.}}}
\end{exoresolu}

\begin{popfigure}
\begin{tikzpicture}[scale=0.7, every node/.style={font=\small}]
% Titre
\node[align=center, draw=poporange, thick, rounded corners, fill=poporange!10, inner sep=8pt] at (0, 5.5) 
{\textbf{Modélisation : Beignets / Boissons}\\ 
Variables : $x$ = nb beignets, $y$ = nb boissons\\
Contraintes : 
\begin{tabular}{rcl}
$150x + 200y$ &=& $1500$ (Budget)\\
$x + y$ &=& $9$ (Quantité)
\end{tabular}
};

% Représentation graphique des droites
\begin{scope}[shift={(0,0)}, scale=0.5]
\draw[->] (-1,0) -- (11,0) node[right] {$x$ (beignets)};
\draw[->] (0,-1) -- (0,11) node[above] {$y$ (boissons)};
% Droite x+y=9
\draw[popblue, thick, domain=-1:10, samples=2] plot (\x, {9-\x}) node[right, pos=0.8, black] {$x+y=9$};
% Droite 150x+200y=1500 -> 3x+4y=30 -> y = 7.5 - 0.75x
\draw[popgreen, thick, domain=-1:10, samples=2] plot (\x, {7.5 - 0.75*\x}) node[right, pos=0.2, black] {$150x+200y=1500$};
% Point intersection (6,3)
\fill[popred] (6,3) circle (4pt) node[above right] {$(6;3)$};
% Grille légère
\foreach \i in {1,...,10} {
    \draw[popmuted, very thin] (\i,-0.1) -- (\i,10.1);
    \draw[popmuted, very thin] (-0.1,\i) -- (10.1,\i);
}
\end{scope}

% Boîte solution finale
\node[align=center, draw=popgreen, thick, rounded corners, fill=popgreen!10, inner sep=8pt] at (0, -3.5) 
{\textbf{Solution : $(x;y) = (6;3)$}\\
6 beignets, 3 boissons\\
Vérif : $6+3=9$ et $900+600=1500$};
\end{tikzpicture}
\legende{Représentation graphique du système : l'intersection des deux droites donne la solution unique $(6;3)$.}
\end{popfigure}

\courssub{Autres familles de situations concrètes}

La structure « Quantité totale / Valeur totale » (ou « Nombre / Prix ») se retrouve dans de nombreux contextes. Voici trois autres cas fréquents au programme de 3ème.

\textbf{Cas 1 : Les mélanges (Concentration / Composition)}\par
On mélange deux produits de compositions différentes pour obtenir un produit intermédiaire.
\textit{Exemple :} On mélange du café « Robusta » à 2500 FCFA/kg et du café « Arabica » à 4000 FCFA/kg pour obtenir 10 kg de mélange vendu 3200 FCFA/kg.
Inconnues : $x$ = masse de Robusta (kg), $y$ = masse d'Arabica (kg).
Système :
\[
\begin{cases}
x + y = 10 & \text{(masse totale)}\\
2500x + 4000y = 3200 \times 10 & \text{(valeur totale = prix final $\times$ masse totale)}
\end{cases}
\]

\textbf{Cas 2 : Géométrie (Périmètre / Aire / Angles)}\par
Les relations géométriques fournissent des équations linéaires.
\textit{Exemple (Rectangle) :} Le périmètre d'un rectangle est 40 cm. Sa longueur dépasse sa largeur de 6 cm. Trouver les dimensions.
Inconnues : $L$ (longueur), $l$ (largeur) en cm.
Système :
\[
\begin{cases}
2(L + l) = 40 \iff L + l = 20\\
L - l = 6
\end{cases}
\]
\emph{Astuce :} Ici, l'\textbf{élimination par addition} est immédiate : $(L+l)+(L-l)=20+6 \Rightarrow 2L=26 \Rightarrow L=13$, puis $l=7$.

\textit{Exemple (Triangle) :} Dans un triangle, un angle vaut $50^{\circ}$. La différence des deux autres est de $20^{\circ}$. Trouver les angles.
Inconnues : $x, y$ (les deux angles inconnus).
Système :
\[
\begin{cases}
x + y + 50 = 180 \iff x + y = 130\\
x - y = 20 \quad (\text{on suppose } x > y)
\end{cases}
\]

\textbf{Cas 3 : Tarification / Transport (Forfait + Unitaire)}\par
Situation très courante (taxi, électricité CAMTEL/ENEO, forfaits internet MTN/Orange).
\textit{Exemple :} Un taxi compteur facture une prise en charge (forfait) + un prix au kilomètre. Un trajet de 5 km coûte 2500 FCFA. Un trajet de 12 km coûte 4600 FCFA. Trouver la prise en charge et le prix du km.
Inconnues : $a$ = prise en charge (FCFA), $b$ = prix du km (FCFA/km).
Système :
\[
\begin{cases}
a + 5b = 2500\\
a + 12b = 4600
\end{cases}
\]
\emph{Méthode :} Élimination par soustraction : $(L_2) - (L_1) \Rightarrow 7b = 2100 \Rightarrow b = 300$. Puis $a = 2500 - 1500 = 1000$.

\begin{savaistu}[La modélisation, cœur de la Recherche Opérationnelle]
Ce que tu fais ici — transformer un problème concret en système d'équations linéaires — est la base de la \textbf{Recherche Opérationnelle (RO)}. Cette discipline mathématique, née pendant la Seconde Guerre mondiale pour optimiser la logistique militaire (ravitaillement, déplacements de troupes), est aujourd'hui partout :
\begin{itemize}
    \item \textbf{Transport/Logistique :} Optimiser les tournées des camions (SOTRA, CAMPOST) pour minimiser le carburant.
    \item \textbf{Télécoms (MTN/Orange) :} Routage des appels et données, dimensionnement des antennes relais.
    \item \textbf{Agriculture/Industrie :} Déterminer le mélange optimal d'engrais ou de matières premières pour maximiser le rendement ou le profit (Programmation Linéaire = systèmes d'inéquations, vu en Première).
    \item \textbf{Finance :} Gestion de portefeuille, équilibre offre/demande sur les marchés.
\end{itemize}
Maîtriser la modélisation en 3ème, c'est poser la première pierre d'un outil de décision puissant utilisé par les ingénieurs, économistes et data scientists du Cameroun et d'ailleurs.
\end{savaistu}

\courssub{Exercices d'entraînement progressifs}

\begin{exercice}[Exercice 1 : Classique « Stylos et Cahiers » (Type BEPC)]
Un commerçant a en stock 20 articles composés de stylos et de cahiers. Le prix d'un stylo est de 300 FCFA, celui d'un cahier est de 500 FCFA. La valeur totale du stock est de 8000 FCFA.
\begin{enumerate}
    \item En posant $x$ le nombre de stylos et $y$ le nombre de cahiers, écrire le système d'équations modélisant la situation.
    \item Résoudre ce système.
    \item Combien y a-t-il de stylos et de cahiers dans le stock ?
\end{enumerate}
\end{exercice}

\begin{corrige}[Exercice 1]
\begin{enumerate}
    \item \textbf{Modélisation.}
    $x$ = nombre de stylos, $y$ = nombre de cahiers. ($x, y \in \mathbb{N}$)
    \begin{itemize}
        \item Quantité : $x + y = 20$
        \item Valeur : $300x + 500y = 8000$
    \end{itemize}
    Système : $\begin{cases} x + y = 20 \\ 300x + 500y = 8000 \end{cases}$ (On peut simplifier la 2ème par 100 : $3x + 5y = 80$).

    \item \textbf{Résolution (Substitution).}
    De la première : $x = 20 - y$.
    Dans la seconde simplifiée : $3(20 - y) + 5y = 80 \iff 60 - 3y + 5y = 80 \iff 2y = 20 \iff y = 10$.
    Alors $x = 20 - 10 = 10$.
    Solution : $(10; 10)$.

    \item \textbf{Interprétation et Réponse.}
    $x=10, y=10$ sont des entiers positifs. Vérification : $10+10=20$ articles. $10\times300 + 10\times500 = 3000+5000=8000$ FCFA.
    \fbox{\textbf{Il y a 10 stylos et 10 cahiers dans le stock.}}
\end{enumerate}
\end{corrige}

\begin{exercice}[Exercice 2 : Problème géométrique (Rectangle)]
Un terrain rectangulaire a un périmètre de 80 mètres. Sa longueur est supérieure de 12 mètres à sa largeur.
\begin{enumerate}
    \item Choisir les inconnues adaptées et écrire le système.
    \item Résoudre par la méthode d'élimination (addition/soustraction).
    \item Calculer l'aire de ce terrain.
\end{enumerate}
\end{exercice}

\begin{corrige}[Exercice 2]
\begin{enumerate}
    \item Soit $L$ la longueur et $l$ la largeur (en mètres). $L, l > 0$.
    \begin{itemize}
        \item Périmètre : $2(L+l) = 80 \iff L + l = 40$.
        \item Différence : $L - l = 12$.
    \end{itemize}
    Système : $\begin{cases} L + l = 40 \\ L - l = 12 \end{cases}$

    \item \textbf{Élimination par addition.}
    On additionne membre à membre les deux équations :
    $(L+l) + (L-l) = 40 + 12 \iff 2L = 52 \iff L = 26$.
    On soustrait la 2ème de la 1ère (ou on remplace) :
    $(L+l) - (L-l) = 40 - 12 \iff 2l = 28 \iff l = 14$.
    (Vérif : $26+14=40$, $26-14=12$). Solution : $(26; 14)$.

    \item \textbf{Aire.}
    $\mathcal{A} = L \times l = 26 \times 14 = 364$ m$^2$.
    \fbox{\textbf{Le terrain mesure 26 m sur 14 m. Son aire est de 364 m$^2$.}}
\end{enumerate}
\end{corrige}

\begin{exercice}[Exercice 3 : Tarification « Forfait + Kilométrique » (Contexte Transport)]
Pour un trajet en taxi « compteur » à Yaoundé, le prix à payer $P$ (en FCFA) s'écrit $P = a + b \times d$, où $d$ est la distance en km, $a$ la prise en charge (forfait) et $b$ le prix du kilomètre.
On sait que :
\begin{itemize}
    \item Un trajet de 4 km coûte 1800 FCFA.
    \item Un trajet de 10 km coûte 3300 FCFA.
\end{itemize}
\begin{enumerate}
    \item Écrire le système d'équations vérifié par $a$ et $b$.
    \item Déterminer $a$ et $b$.
    \item Combien coûtera un trajet de 15 km ?
\end{enumerate}
\end{exercice}

\begin{corrige}[Exercice 3]
\begin{enumerate}
    \item Traduction :
    \begin{align*}
    4\text{ km} &: a + 4b = 1800 \quad (1)\\
    10\text{ km} &: a + 10b = 3300 \quad (2)
    \end{align*}
    Système : $\begin{cases} a + 4b = 1800 \\ a + 10b = 3300 \end{cases}$

    \item \textbf{Élimination par soustraction $(2) - (1)$ :}
    $(a + 10b) - (a + 4b) = 3300 - 1800 \iff 6b = 1500 \iff b = 250$.
    Le prix du kilomètre est de 250 FCFA/km.
    On trouve $a$ avec (1) : $a + 4(250) = 1800 \iff a + 1000 = 1800 \iff a = 800$.
    La prise en charge est de 800 FCFA.

    \item \textbf{Prédiction pour 15 km :}
    $P = a + 15b = 800 + 15 \times 250 = 800 + 3750 = 4550$ FCFA.
    \fbox{\textbf{La prise en charge est 800 FCFA, le km coûte 250 FCFA. Un trajet de 15 km coûte 4550 FCFA.}}
\end{enumerate}
\end{corrige}

\courssub{Synthèse : La check-list du modélisateur expert}

Avant de rendre ta copie au BEPC (ou dans la vie), passe ce check-list mental :

\begin{center}
\begin{tabularx}{\linewidth}{|l|X|}\hline
\textbf{Étape} & \textbf{Question à se poser} \\ \hline
1. Variables & Ai-je écrit \og Soit $x$ ... et $y$ ... \fg avec les unités ? \\ \hline
2. Système & Mes deux équations traduisent-elles bien \emph{deux infos distinctes} du texte ? \\ \hline
3. Résolution & Ai-je choisi la méthode la plus simple (substitution / élimination / Cramer) ? Mes calculs sont-ils posés proprement ? \\ \hline
4. Cohérence & La solution $(x;y)$ est-elle dans le bon ensemble ($\mathbb{N}$, $\mathbb{R}_+$, $]0;180[$ pour un angle...) ? \\ \hline
5. Vérification & Ai-je remis $(x;y)$ dans les \emph{deux} équations initiales (ou dans le texte) ? \\ \hline
6. Phrase finale & Ma réponse est-elle une phrase complète, en français, avec les unités, qui répond \emph{exactement} à la question posée ? \\ \hline
\end{tabularx}
\end{center}

\begin{aretenir}[L'essentiel à retenir pour le BEPC]
\begin{itemize}
    \item Un problème à deux inconnues se modélise par un \textbf{système de deux équations du premier degré}.
    \item Les mots-clés pour repérer les équations : \og total \fg, \og somme \fg, \og différence \fg, \og produit \fg (prix $\times$ quantité), \og périmètre \fg, \og somme des angles \fg.
    \item La \textbf{résolution} n'est qu'une étape intermédiaire. L'étape cruciale est l'\textbf{interprétation} : la solution mathématique doit avoir un sens physique (entier, positif, plausible).
    \item La \textbf{rédaction finale} (phrase de réponse) fait partie de la note. Ne l'oublie jamais.
\end{itemize}
\end{aretenir}

\begin{savaistu}[Erreur de modélisation célèbre : L'avion et le vent]
Un grand classique des concours (et de la vie réelle pour les pilotes) : \og Un avion fait l'aller-retour entre Yaoundé et Douala (distance $D$). Sa vitesse propre est $v$. Le vent souffle dans le sens Yaoundé-Douala à la vitesse $w$. Exprimer le temps total. \fg
Beaucoup d'élèves écrivent : $t = \frac{D}{v+w} + \frac{D}{v+w}$ (oubliant que le vent est \emph{contre} au retour).
Le bon système/modèle : $t = \frac{D}{v+w} + \frac{D}{v-w}$.
La modélisation, c'est \textbf{respecter la physique du problème}, pas seulement aligner des symboles.
\end{savaistu}

\newpage

\coursec{Activité d'intégration}

\coursec{Équations du premier degré dans $\mathbb{R}\times\mathbb{R}$ : Modélisation et résolution — Activité d'intégration}

\courssub{Objectifs de l'activité}

À l'issue de cette activité d'intégration, tu seras capable de :
\begin{itemize}
\item \cle{modéliser} une situation concrète de la vie courante (gestion d'une cantine scolaire) par un système de deux équations du premier degré à deux inconnues ;
\item choisir de façon pertinente les \cle{inconnues} et traduire chaque contrainte de l'énoncé en équation ou en inéquation ;
\item résoudre un système par la méthode de \cle{substitution} ou d'\cle{élimination} (combinaison linéaire), puis vérifier le résultat par le \cle{déterminant} ;
\item discuter le \cle{nombre de solutions} d'un système (solution unique, aucune solution, infinité de solutions) selon les données ;
\item \cle{rédiger et justifier} une démarche complète, une réponse et une conclusion, puis la présenter oralement ;
\item prendre une \cle{décision budgétaire} argumentée, en proposant une solution entière réaliste.
\end{itemize}

\courssub{Situation}

\textbf{La cantine du Collège Bilingue de Nkol-Eton.}\par
Le Collège Bilingue de Nkol-Eton, situé dans la région du Centre, organise chaque jour une cantine pour ses élèves. La présidente de la coopérative scolaire, Mme Abena, est chargée de planifier les repas de la journée de clôture de l'année.

Voici le cahier des charges qu'elle a reçu du directeur :
\begin{itemize}
\item la cantine doit préparer exactement \textbf{120 repas}, chaque repas étant composé d'un \textbf{plat principal} (riz sauté au poulet) et éventuellement d'un \textbf{dessert} (salade de fruits) ;
\item le \textbf{budget total} disponible est de \textbf{180\,000 FCFA} ;
\item un plat principal coûte \textbf{1\,200 FCFA} et un dessert coûte \textbf{300 FCFA} ;
\item pour des raisons de diététique, le nombre de desserts servis doit être \textbf{au moins le double} du nombre de plats principaux... mais attention : chaque repas contient un plat, donc il y a 120 plats ! Mme Abena se rend compte que cette contrainte est impossible et demande au directeur de la corriger.
\end{itemize}

Après discussion, le cahier des charges définitif est le suivant :

\begin{aretenir}[Cahier des charges définitif]
\begin{itemize}
\item Nombre total d'éléments à préparer (plats principaux $+$ desserts) : \textbf{240 éléments} pour 120 élèves (chaque élève reçoit un plat et un dessert, et on prépare des éléments supplémentaires pour les enseignants et les invités).
\item Budget total : \textbf{180\,000 FCFA}, à dépenser \textbf{exactement}.
\item Coût d'un plat : \textbf{1\,200 FCFA} ; coût d'un dessert : \textbf{300 FCFA}.
\item Contrainte diététique : le nombre de desserts doit être \textbf{au moins égal au nombre de plats}.
\end{itemize}
\end{aretenir}

Mme Abena, qui a été une excellente élève en mathématiques, se pose les questions suivantes : combien de plats et combien de desserts doit-elle commander pour dépenser exactement le budget ? La contrainte diététique est-elle respectée ? Et si le budget avait été différent, le problème aurait-il encore une solution ?

\courssub{Consignes}

Travail en groupes de 4 élèves. Chaque groupe rédige sa démarche complète sur une feuille, puis un rapporteur présente les résultats au tableau.

\begin{enumerate}
\item \textbf{Choix des inconnues.} Désigner par deux lettres les deux quantités inconnues du problème. Préciser clairement ce que chaque lettre représente.
\item \textbf{Mise en équations.} Traduire chacune des contraintes suivantes par une équation ou une inéquation :
\begin{enumerate}
\item le nombre total d'éléments à préparer ;
\item la dépense totale exacte de 180\,000 FCFA ;
\item la contrainte diététique (desserts au moins aussi nombreux que les plats).
\end{enumerate}
\item \textbf{Résolution algébrique.} Résoudre le système de deux équations obtenu :
\begin{enumerate}
\item par la méthode de \textbf{substitution} ;
\item par la méthode d'\textbf{élimination} (combinaison linéaire) ;
\item vérifier que les deux méthodes donnent le même couple solution.
\end{enumerate}
\item \textbf{Étude du déterminant.} Calculer le déterminant du système. Que permet-il d'affirmer sur le nombre de solutions avant même de résoudre ?
\item \textbf{Vérification des contraintes.} Vérifier que la solution trouvée respecte la contrainte diététique et le budget. Les valeurs sont-elles réalistes (nombres entiers positifs) ?
\item \textbf{Discussion.} Répondre aux questions suivantes en justifiant :
\begin{enumerate}
\item Si le budget avait été de \textbf{150\,000 FCFA} au lieu de 180\,000 FCFA, le système aurait-il encore une solution entière positive ? Que peut-on en conclure ?
\item Si l'on \textbf{relâche} la contrainte d'égalité du budget (on accepte de dépenser \emph{au plus} 180\,000 FCFA), combien de solutions entières peut-on envisager ? Donner deux exemples.
\end{enumerate}
\item \textbf{Conclusion et décision.} Rédiger un court paragraphe de conseil à Mme Abena : que doit-elle commander, et pourquoi ?
\end{enumerate}

\courssub{Ressources à mobiliser}

\begin{aretenir}[Boîte à outils de la leçon]
\begin{itemize}
\item Une \cle{équation du premier degré à deux inconnues} est une équation de la forme $ax + by = c$ avec $(a\,;\,b) \neq (0\,;\,0)$. Ses solutions sont des couples $(x\,;\,y)$, représentés graphiquement par les points d'une \textbf{droite}.
\item Un \cle{système de deux équations} à deux inconnues s'écrit $\begin{cases} ax + by = c \\ a'x + b'y = c' \end{cases}$. Le résoudre, c'est trouver tous les couples $(x\,;\,y)$ vérifiant les deux équations à la fois.
\item \textbf{Méthode de substitution} : exprimer une inconnue en fonction de l'autre dans une équation, puis remplacer dans l'autre équation.
\item \textbf{Méthode d'élimination} : multiplier les équations par des nombres bien choisis puis les additionner membre à membre pour éliminer une inconnue.
\item Le \cle{déterminant} du système est $\Delta = ab' - a'b$ :
\begin{itemize}
\item si $\Delta \neq 0$ : le système admet une \textbf{unique solution} ;
\item si $\Delta = 0$ : le système admet \textbf{aucune solution} ou une \textbf{infinité de solutions} (droites parallèles distinctes ou confondues).
\end{itemize}
\item Une solution d'un problème concret doit être \textbf{vérifiée} dans l'énoncé et être \textbf{compatible avec le contexte} (ici : nombres entiers positifs).
\end{itemize}
\end{aretenir}

\courssub{Démarche et corrigé commenté}

\begin{exoresolu}[La cantine du Collège Bilingue de Nkol-Eton — corrigé complet]
\textbf{Énoncé rappelé.} Préparer 240 éléments (plats à 1\,200 FCFA et desserts à 300 FCFA) en dépensant exactement 180\,000 FCFA, avec au moins autant de desserts que de plats. Déterminer le nombre de plats et de desserts, puis discuter.
\tcblower
\textbf{Étape 1 : choix des inconnues.}\par
Soit $x$ le nombre de \textbf{plats principaux} et $y$ le nombre de \textbf{desserts}. Ce sont des nombres entiers naturels.

\textbf{Étape 2 : mise en équations.}\par
Le nombre total d'éléments est 240, donc : $x + y = 240$.\par
La dépense totale est de 180\,000 FCFA : $x$ plats à 1\,200 FCFA et $y$ desserts à 300 FCFA coûtent $1\,200x + 300y$ francs, d'où : $1\,200x + 300y = 180\,000$.\par
On peut simplifier cette équation en divisant par 300 : $4x + y = 600$.\par
La contrainte diététique s'écrit : $y \geq x$.\par
Le système à résoudre est donc :
\[ (S) : \begin{cases} x + y = 240 \\ 4x + y = 600 \end{cases} \qquad \text{avec la contrainte } y \geq x. \]

\textbf{Étape 3 : résolution par substitution.}\par
De la première équation, on tire $y = 240 - x$. On remplace dans la deuxième :
\begin{align*}
4x + (240 - x) &= 600 \\
3x + 240 &= 600 \\
3x &= 360 \\
x &= 120.
\end{align*}
Alors $y = 240 - 120 = 120$.\par
Le couple solution est $(120\,;\,120)$.

\textbf{Étape 4 : vérification par élimination.}\par
On soustrait la première équation de la deuxième, membre à membre :
\[ (4x + y) - (x + y) = 600 - 240 \quad\Longrightarrow\quad 3x = 360 \quad\Longrightarrow\quad x = 120. \]
En reportant dans $x + y = 240$ : $y = 120$. On retrouve bien le même couple $(120\,;\,120)$ : les deux méthodes concordent.

\textbf{Étape 5 : étude du déterminant.}\par
Avec $a = 1$, $b = 1$, $a' = 4$, $b' = 1$ :
\[ \Delta = ab' - a'b = 1 \times 1 - 4 \times 1 = 1 - 4 = -3. \]
Comme $\Delta = -3 \neq 0$, le système admet une \textbf{unique solution}. Cela confirme que le couple $(120\,;\,120)$ est le seul possible. Graphiquement, les deux droites d'équations $x + y = 240$ et $4x + y = 600$ sont sécantes en un seul point.

\begin{popfigure}
\begin{tikzpicture}
\begin{axis}[
    width=\linewidth,
    height=0.6\linewidth,
    axis lines=middle,
    xlabel=$x$ (plats), ylabel=$y$ (desserts),
    xmin=0, xmax=250,
    ymin=0, ymax=250,
    xtick={0,60,120,180,240},
    ytick={0,60,120,180,240},
    grid=major,
    grid style={popline},
    tick label style={font=\small},
    label style={font=\small},
    legend style={at={(0.98,0.98)}, anchor=north east, font=\small, fill=popcream, draw=popmuted},
]
\addplot[popblue, thick, domain=0:240, samples=2] {240-x};
\addlegendentry{$x+y=240$}
\addplot[poporange, thick, domain=0:150, samples=2] {600-4*x};
\addlegendentry{$4x+y=600$}
\addplot[only marks, mark=*, mark size=3, popgreen] coordinates {(120,120)};
\addlegendentry{Solution $(120;120)$}
\end{axis}
\end{tikzpicture}
\legende{Représentation graphique du système : intersection des deux droites en $(120;120)$.}
\end{popfigure}

\textbf{Étape 6 : vérification des contraintes concrètes.}
\begin{itemize}
\item Total : $120 + 120 = 240$ éléments. \textbf{Correct.}
\item Dépense : $1\,200 \times 120 + 300 \times 120 = 144\,000 + 36\,000 = 180\,000$ FCFA. \textbf{Budget exactement dépensé.}
\item Contrainte diététique : $y = 120 \geq x = 120$. \textbf{Respectée} (il y a autant de desserts que de plats).
\item Les valeurs $120$ et $120$ sont des entiers positifs : la solution est \textbf{réaliste}.
\end{itemize}

\textbf{Étape 7 : discussion.}\par
\emph{(a) Budget de 150\,000 FCFA.} Le système devient :
\[ \begin{cases} x + y = 240 \\ 4x + y = 500 \end{cases} \quad\text{(après division par 300).} \]
Par élimination : $3x = 260$, donc $x = \dfrac{260}{3} \approx 86,7$, qui n'est \textbf{pas un entier}. Le déterminant vaut toujours $-3 \neq 0$, donc il existe une solution mathématique unique, mais elle n'est \textbf{pas réaliste} : on ne peut pas commander $86,7$ plats ! De plus $y = 240 - \dfrac{260}{3} = \dfrac{460}{3} \approx 153,3$. Conclusion : avec 150\,000 FCFA, Mme Abena ne peut pas dépenser exactement son budget ; elle devra dépenser \emph{moins} (par exemple 86 plats et 154 desserts coûtent $103\,200 + 46\,200 = 149\,400$ FCFA).\par
\emph{(b) Budget relâché (au plus 180\,000 FCFA).} On a alors le système mixte :
\[ \begin{cases} x + y = 240 \\ 4x + y \leq 600 \end{cases} \quad\Longleftrightarrow\quad 3x \leq 360 \quad\Longleftrightarrow\quad x \leq 120. \]
Il existe alors une \textbf{infinité de solutions} mathématiques, et de nombreuses solutions entières : par exemple $(x\,;\,y) = (100\,;\,140)$ qui coûte $120\,000 + 42\,000 = 162\,000$ FCFA, ou $(x\,;\,y) = (110\,;\,130)$ qui coûte $132\,000 + 39\,000 = 171\,000$ FCFA. Toutes respectent $y \geq x$.

\textbf{Étape 8 : conclusion et décision.}\par
Mme Abena doit commander \textbf{120 plats principaux et 120 desserts}, pour une dépense totale exacte de 180\,000 FCFA. Cette solution est l'unique solution du système, elle respecte la contrainte diététique et épuise exactement le budget. Si un imprévu réduisait le budget, elle devrait diminuer le nombre de plats (qui sont les éléments les plus coûteux) tout en gardant $x + y = 240$.
\end{exoresolu}

\begin{attention}[Erreurs fréquentes à éviter]
\begin{itemize}
\item \textbf{Oublier de définir les inconnues} : écrire « soit $x$ et $y$ » sans préciser ce qu'ils représentent est une erreur de rédaction sanctionnée au BEPC.
\item \textbf{Se tromper d'équation} : le total des éléments donne $x + y = 240$, pas $x + y = 120$ ; relire l'énoncé avant de modéliser.
\item \textbf{Oublier la vérification} : un couple solution doit toujours être testé dans \emph{les deux} équations et dans les contraintes concrètes (entiers positifs, inégalité).
\item \textbf{Conclure trop vite} : si $\Delta = 0$, on ne peut pas conclure sans examiner si les droites sont confondues (infinité de solutions) ou strictement parallèles (aucune solution).
\item \textbf{Accepter une solution non entière} dans un problème de dénombrement : $\dfrac{260}{3}$ plats n'a aucun sens concret ; il faut le signaler et proposer une solution entière proche.
\end{itemize}
\end{attention}

\courssub{Bilan de l'activité}

Cette activité d'intégration a permis de mobiliser l'ensemble des savoirs et savoir-faire de la leçon dans une situation authentique de gestion budgétaire :

\begin{itemize}
\item \textbf{Modélisation} : le passage d'un cahier des charges rédigé en français à un système d'équations est l'étape la plus délicate ; elle exige de bien choisir les inconnues et de traduire chaque contrainte séparément.
\item \textbf{Résolution algébrique} : les deux méthodes (substitution et élimination) ont conduit au même couple $(120\,;\,120)$, ce qui illustre leur équivalence ; le déterminant $\Delta = -3 \neq 0$ a permis d'affirmer \emph{a priori} l'unicité de la solution.
\item \textbf{Interprétation graphique} : résoudre le système revient à chercher le point d'intersection de deux droites ; le signe du déterminant traduit la position relative de ces droites.
\item \textbf{Discussion du nombre de solutions} : en modifiant le budget, on a rencontré un cas sans solution entière réaliste ; en relâchant la contrainte d'égalité, on est passé d'une solution unique à une infinité de solutions. C'est le cœur de la leçon : le nombre de solutions dépend des données.
\item \textbf{Prise de décision argumentée} : la mathématique ne s'arrête pas au calcul ; il faut vérifier la cohérence concrète (entiers positifs, contrainte diététique) et rédiger une conclusion utile à la décision.
\end{itemize}

\begin{savaistu}[Les systèmes d'équations, un outil de gestion quotidien]
Les coopératives scolaires, les comités d'organisation de la Fête de la Jeunesse ou les gestionnaires de transport urbain (comme les agences de voyage de la gare routière de Mvan à Yaoundé) résolvent chaque semaine, souvent sans le savoir, des systèmes d'équations du premier degré : combien de tickets de deux catégories vendre pour atteindre une recette donnée, comment répartir un budget entre plusieurs postes de dépenses... Au-delà du BEPC, ces problèmes de répartition optimale sont au cœur d'un domaine des mathématiques appelé \emph{recherche opérationnelle}, utilisé par les grandes entreprises de télécommunications comme MTN et Orange pour planifier leurs réseaux au Cameroun.
\end{savaistu}

\newpage

\coursec{Méthodes et exercices résolus}

\coursec{Leçon 11 : Équations du premier degré dans $\mathbb{R}\times\mathbb{R}$}

\courssub{Équation du premier degré à deux inconnues : vocabulaire et représentation}

\begin{definition}[Équation du premier degré à deux inconnues]
Une \cle{équation du premier degré à deux inconnues} $x$ et $y$ s'écrit sous la forme canonique :
\[ ax + by = c \]
où $a$, $b$, $c$ sont des nombres réels donnés, avec $a$ et $b$ non tous les deux nuls.
Un \cle{couple ordonné} $(x_0\,;y_0)$ est une \cle{solution} de l'équation si l'égalité $ax_0 + by_0 = c$ est vraie.
L'ensemble des solutions se note $\mathcal{S}$.
\end{definition}

\begin{definition}[Système de deux équations du premier degré]
Un \cle{système de deux équations du premier degré} à deux inconnues $x$ et $y$ est un ensemble de deux équations :
\[ \begin{cases} a_1x + b_1y = c_1 \\ a_2x + b_2y = c_2 \end{cases} \]
Un couple $(x_0\,;y_0)$ est \cle{solution du système} s'il vérifie \emph{les deux équations simultanément}.
L'ensemble des solutions du système se note $\mathcal{S}$.
\end{definition}

\begin{propriete}[Représentation graphique]
Dans un repère orthonormé $(O\,;\vec{\imath}\,;\vec{\jmath})$, l'ensemble des solutions d'une équation $ax+by=c$ (avec $(a,b)\neq(0,0)$) est une \cle{droite}.
\begin{itemize}
\item Si $b\neq 0$, la droite a pour équation réduite $y = -\dfrac{a}{b}x + \dfrac{c}{b}$ (coefficient directeur $-\frac{a}{b}$, ordonnée à l'origine $\frac{c}{b}$).
\item Si $b=0$, l'équation est $ax=c$, soit $x=\frac{c}{a}$ : c'est une droite parallèle à l'axe des ordonnées.
\end{itemize}
Résoudre un système de deux équations, c'est chercher les coordonnées du(des) point(s) d'intersection des deux droites associées.
\end{propriete}

\begin{methode}[Vérifier qu'un couple $(x\,;y)$ est solution d'une équation ou d'un système]
\begin{enumerate}
\item \textbf{Repérer} l'inconnue $x$ (premier nombre) et l'inconnue $y$ (second nombre) dans le couple proposé.
\item \textbf{Remplacer} $x$ et $y$ par leurs valeurs dans le(s) membre(s) de gauche.
\item \textbf{Calculer} le membre de gauche et le comparer au membre de droite.
\item Pour un \textbf{système}, le couple doit vérifier \cle{les deux équations}. S'il en échoue une seule, il n'est pas solution.
\item \textbf{Conclure} par une phrase : « Le couple $(x_0\,;y_0)$ est (ou n'est pas) solution de... ».
\end{enumerate}
\end{methode}

\begin{exoresolu}[Vérification d'un couple solution]
On considère le système $\begin{cases} 2x + 3y = 1 \\ x - 4y = 6 \end{cases}$.
Le couple $(2\,;-1)$ est-il solution ? Le couple $(1\,;1)$ l'est-il ?
\tcblower
\textbf{Test du couple $(2\,;-1)$ :} $x=2$, $y=-1$.
\begin{itemize}
\item Première équation : $2\times 2 + 3\times(-1) = 4 - 3 = 1$ (membre de droite $=1$) \checkmark
\item Deuxième équation : $2 - 4\times(-1) = 2 + 4 = 6$ (membre de droite $=6$) \checkmark
\end{itemize}
Le couple vérifie les deux équations : $(2\,;-1)$ \textbf{est solution} du système.

\textbf{Test du couple $(1\,;1)$ :} $x=1$, $y=1$.
\begin{itemize}
\item Première équation : $2\times 1 + 3\times 1 = 5 \neq 1$.
\end{itemize}
Inutile de tester la seconde : $(1\,;1)$ \textbf{n'est pas solution} du système.
\end{exoresolu}

\courssub{Système de deux équations : approche graphique}

\begin{methode}[Résoudre graphiquement un système]
\begin{enumerate}
\item \textbf{Mettre chaque équation sous forme réduite} $y = ax + b$ (ou $x = \text{constante}$) en isolant $y$.
\item \textbf{Tracer chaque droite} dans un repère orthonormé : calculer deux points par droite (ex. $x=0$ et $x=1$ ou $x=0$ et $y=0$), les placer, tracer à la règle.
\item \textbf{Lire les coordonnées du point d'intersection} : c'est la solution $(x\,;y)$.
\item \textbf{Vérifier algébriquement} le couple lu dans les deux équations initiales (la lecture graphique est approximative).
\item \textbf{Cas particuliers :}
\begin{itemize}
\item Droites \textbf{strictement parallèles} (même coefficient directeur, ordonnées à l'origine différentes) $\Rightarrow$ \cle{aucune solution} ($\mathcal{S}=\varnothing$).
\item Droites \textbf{confondues} $\Rightarrow$ \cle{une infinité de solutions} (tous les points de la droite).
\end{itemize}
\end{enumerate}
\end{methode}

\begin{exoresolu}[Résolution graphique]
Résoudre graphiquement le système $\begin{cases} x + y = 5 \\ 2x - y = 1 \end{cases}$, puis vérifier par le calcul.
\tcblower
\textbf{Étape 1 : Forme réduite.}
$y = -x + 5$ \quad et \quad $y = 2x - 1$.

\textbf{Étape 2 : Points de passage.}
Pour $y=-x+5$ : $A(0\,;5)$ et $B(5\,;0)$.
Pour $y=2x-1$ : $C(0\,;-1)$ et $D(2\,;3)$.

\textbf{Étape 3 : Tracé et lecture.}
\begin{popfigure}
\begin{tikzpicture}[scale=0.6]
\draw[popline,->] (-1,0) -- (6.5,0) node[below]{$x$};
\draw[popline,->] (0,-2) -- (0,6) node[left]{$y$};
\foreach \i in {1,...,6} \draw[popline] (\i,0.1) -- (\i,-0.1) node[below,font=\scriptsize]{\i};
\foreach \j in {-1,...,5} \draw[popline] (0.1,\j) -- (-0.1,\j) node[left,font=\scriptsize]{\j};
\draw[popblue,thick,domain=-0.5:5.5] plot (\x,{-\x+5}) node[right,font=\small,popblue]{$y=-x+5$};
\draw[poporange,thick,domain=-0.2:3.2] plot (\x,{2*\x-1}) node[above right,font=\small,poporange]{$y=2x-1$};
\fill[popink] (2,3) circle (3pt) node[above right,font=\small]{$(2\,;3)$};
\draw[popmuted,dashed] (2,0) node[below,font=\scriptsize]{$2$} -- (2,3) -- (0,3) node[left,font=\scriptsize]{$3$};
\end{tikzpicture}
\legende{Les droites se coupent au point $I(2\,;3)$.}
\end{popfigure}
Les droites se coupent au point $I(2\,;3)$. Solution graphique : $(2\,;3)$.

\textbf{Étape 4 : Vérification.}
$2 + 3 = 5$ \checkmark \quad et \quad $2\times 2 - 3 = 1$ \checkmark.

\textbf{Conclusion :} Le système admet pour unique solution $\boxed{(2\,;3)}$.
\end{exoresolu}

\courssub{Méthodes algébriques : substitution et élimination (combinaison linéaire)}

\begin{methode}[Méthode de substitution]
À privilégier quand une inconnue a pour coefficient $1$ ou $-1$ dans l'une des équations.
\begin{enumerate}
\item \textbf{Isoler} une inconnue dans l'équation la plus simple (ex. $x = 5 - y$).
\item \textbf{Substituer} cette expression dans l'autre équation : on obtient une équation à une inconnue.
\item \textbf{Résoudre} cette équation.
\item \textbf{Revenir} à l'expression de l'étape 1 pour calculer la seconde inconnue.
\item \textbf{Vérifier} dans les deux équations de départ et conclure : $\mathcal{S} = \{(x\,;y)\}$.
\end{enumerate}
\end{methode}

\begin{exoresolu}[Résolution par substitution]
Résoudre dans $\mathbb{R}\times\mathbb{R}$ : $\begin{cases} x + 2y = 7 \\ 3x - y = 7 \end{cases}$.
\tcblower
\textbf{Étape 1 :} Dans la première, $x$ a pour coefficient $1$ : $x = 7 - 2y$.

\textbf{Étape 2 :} Substitution dans la seconde : $3(7 - 2y) - y = 7$.

\textbf{Étape 3 :} Résolution :
\begin{align*}
21 - 6y - y &= 7 \\
21 - 7y &= 7 \\
-7y &= -14 \\
y &= 2
\end{align*}

\textbf{Étape 4 :} $x = 7 - 2\times 2 = 3$.

\textbf{Étape 5 : Vérification.}
$3 + 2\times 2 = 7$ \checkmark \quad ; \quad $3\times 3 - 2 = 7$ \checkmark.

\textbf{Conclusion :} $\mathcal{S} = \{(3\,;2)\}$.
\end{exoresolu}

\begin{methode}[Méthode d'élimination (combinaison linéaire)]
À privilégier quand aucune inconnue n'a un coefficient $\pm 1$.
\begin{enumerate}
\item \textbf{Choisir} l'inconnue à éliminer et repérer ses coefficients.
\item \textbf{Multiplier} chaque équation par un nombre pour obtenir des coefficients \cle{opposés} devant cette inconnue.
\item \textbf{Additionner} membre à membre les deux équations : l'inconnue disparaît.
\item \textbf{Résoudre} l'équation à une inconnue obtenue.
\item \textbf{Trouver la seconde inconnue} (en remplaçant ou en réitérant l'élimination).
\item \textbf{Vérifier} et conclure.
\end{enumerate}
\end{methode}

\begin{exoresolu}[Résolution par combinaison linéaire]
Résoudre dans $\mathbb{R}\times\mathbb{R}$ : $\begin{cases} 3x + 2y = 16 \quad (E_1)\\ 5x - 4y = 1 \quad (E_2) \end{cases}$.
\tcblower
\textbf{Étape 1 :} On élimine $y$ (coefficients $2$ et $-4$).

\textbf{Étape 2 :} On multiplie $(E_1)$ par $2$ : $6x + 4y = 32 \quad (E_1')$.

\textbf{Étape 3 :} Addition $(E_1') + (E_2)$ :
\[ (6x + 4y) + (5x - 4y) = 32 + 1 \quad\Rightarrow\quad 11x = 33 \]

\textbf{Étape 4 :} $x = 3$.

\textbf{Étape 5 :} On élimine $x$ pour trouver $y$ directement. On multiplie $(E_1)$ par $5$ et $(E_2)$ par $3$ :
\begin{align*}
5(E_1) &: 15x + 10y = 80 \\
3(E_2) &: 15x - 12y = 3
\end{align*}
Soustraction : $(15x + 10y) - (15x - 12y) = 80 - 3 \;\Rightarrow\; 22y = 77 \;\Rightarrow\; y = \frac{77}{22} = \frac{7}{2} = 3,5$.

\textbf{Vérification :} $3\times 3 + 2\times 3,5 = 9 + 7 = 16$ \checkmark ; $5\times 3 - 4\times 3,5 = 15 - 14 = 1$ \checkmark.

\textbf{Conclusion :} $\mathcal{S} = \{(3\,;3,5)\}$.
\end{exoresolu}

\courssub{Discussion du nombre de solutions : déterminant et cas particuliers}

\begin{methode}[Discussion par comparaison des rapports]
Pour le système $\begin{cases} ax + by = c \\ a'x + b'y = c' \end{cases}$, on compare les rapports $\dfrac{a}{a'}$, $\dfrac{b}{b'}$, $\dfrac{c}{c'}$ (dénominateurs non nuls) :
\begin{enumerate}
\item Si $\dfrac{a}{a'} \neq \dfrac{b}{b'}$ : droites \textbf{sécantes} $\Rightarrow$ \cle{une unique solution}.
\item Si $\dfrac{a}{a'} = \dfrac{b}{b'} \neq \dfrac{c}{c'}$ : droites \textbf{strictement parallèles} $\Rightarrow$ \cle{aucune solution} ($\mathcal{S}=\varnothing$).
\item Si $\dfrac{a}{a'} = \dfrac{b}{b'} = \dfrac{c}{c'}$ : droites \textbf{confondues} $\Rightarrow$ \cle{une infinité de solutions}.
\end{enumerate}
\emph{Astuce :} Mettre sous forme $y = mx + p$ pour comparer directement les coefficients directeurs $m$.
\end{methode}

\begin{propriete}[Déterminant et nombre de solutions]
Pour le système $\begin{cases} ax + by = c \\ a'x + b'y = c' \end{cases}$, on définit le \cle{déterminant} $\Delta = ab' - a'b$.
\begin{itemize}
\item Si $\Delta \neq 0$ : le système admet \textbf{une unique solution} (droites sécantes).
\item Si $\Delta = 0$ : les droites sont parallèles ou confondues.
\begin{itemize}
\item Si de plus $ac' - a'c \neq 0$ (ou $bc' - b'c \neq 0$) : \textbf{aucune solution} (droites strictement parallèles).
\item Si $ac' - a'c = 0$ et $bc' - b'c = 0$ : \textbf{infinité de solutions} (droites confondues, équations proportionnelles).
\end{itemize}
\end{itemize}
\end{propriete}

\begin{exoresolu}[Discussion du nombre de solutions]
Sans les résoudre, déterminer le nombre de solutions de :
\[ (S_1) : \begin{cases} 2x - 3y = 5 \\ 4x + y = 3 \end{cases} \quad
(S_2) : \begin{cases} x - 2y = 4 \\ 3x - 6y = 5 \end{cases} \quad
(S_3) : \begin{cases} x - 2y = 4 \\ 3x - 6y = 12 \end{cases} \]
\tcblower
\textbf{Système $(S_1)$ :} $\dfrac{2}{4} = \dfrac{1}{2}$ et $\dfrac{-3}{1} = -3$. $\frac{1}{2} \neq -3$ $\Rightarrow$ \textbf{unique solution}.
(Determinant $\Delta = 2\times 1 - 4\times(-3) = 2 + 12 = 14 \neq 0$).

\textbf{Système $(S_2)$ :} $\dfrac{1}{3}$, $\dfrac{-2}{-6} = \dfrac{1}{3}$, $\dfrac{4}{5}$. $\frac{1}{3} = \frac{1}{3} \neq \frac{4}{5}$ $\Rightarrow$ \textbf{aucune solution}.
($\Delta = 1\times(-6) - 3\times(-2) = -6 + 6 = 0$ et $1\times 5 - 3\times 4 = -7 \neq 0$).

\textbf{Système $(S_3)$ :} $\dfrac{1}{3}$, $\dfrac{-2}{-6} = \dfrac{1}{3}$, $\dfrac{4}{12} = \dfrac{1}{3}$. Trois rapports égaux $\Rightarrow$ \textbf{infinité de solutions}.
($\Delta = 0$ et la 2ème équation $= 3\times$ la 1ère). $\mathcal{S} = \{(x\,;y) \in \mathbb{R}^2 \mid x - 2y = 4\}$.
\end{exoresolu}

\begin{savaistu}[De la Chine antique au GPS]
La résolution de systèmes linéaires est documentée dans le \emph{Jiuzhang Suanshu} (Neuf chapitres sur l'art mathématique, Chine, ~Ier siècle av. J.-C.) avec la méthode des « tableaux » (ancêtre de l'élimination de Gauss). Aujourd'hui, votre téléphone utilise la résolution de systèmes d'équations (trilatération) pour calculer votre position GPS à partir des signaux de plusieurs satellites. En économie, les modèles d'équilibre général (Leontief, prix Nobel 1973) reposent sur la résolution de systèmes géants ($n$ équations, $n$ inconnues) pour prévoir la production nationale.
\end{savaistu}

\courssub{Modélisation et résolution de problèmes concrets}

\begin{methode}[Mettre un problème en système d'équations]
\begin{enumerate}
\item \textbf{Choisir les inconnues} : nommer clairement $x$ et $y$ avec leur sens et leur unité (ex. « $x$ : prix d'un kg de tomates en FCFA »).
\item \textbf{Traduire} chaque information de l'énoncé en une équation reliant $x$ et $y$.
\item \textbf{Résoudre} le système par substitution ou combinaison (choisir la plus simple).
\item \textbf{Vérifier} la cohérence avec le contexte (valeurs positives, entières si comptage...).
\item \textbf{Conclure} par une phrase répondant exactement à la question, avec les unités.
\end{enumerate}
\end{methode}

\begin{exoresolu}[Problème de la vie courante : prix au marché]
Au marché Mokolo, Maman Aïcha achète $3$ kg de tomates et $2$ kg d'oignons pour $\SI{2900}{FCFA}$. Le lendemain, aux mêmes prix, elle achète $2$ kg de tomates et $5$ kg d'oignons pour $\SI{4050}{FCFA}$.
Déterminer le prix d'un kilogramme de tomates et celui d'un kilogramme d'oignons.
\tcblower
\textbf{1. Inconnues :} Soit $x$ le prix du kg de tomates (en FCFA) et $y$ le prix du kg d'oignons (en FCFA).

\textbf{2. Système :}
\[ \begin{cases} 3x + 2y = 2900 \quad (E_1) \\ 2x + 5y = 4050 \quad (E_2) \end{cases} \]

\textbf{3. Résolution par combinaison (élimination de $y$) :}
On multiplie $(E_1)$ par $5$ et $(E_2)$ par $2$ pour obtenir $10y$ et $10y$ (coefficients égaux, on soustraira).
\begin{align*}
5\times(E_1) &: 15x + 10y = 14500 \quad (E_1') \\
2\times(E_2) &: 4x + 10y = 8100 \quad (E_2')
\end{align*}
$(E_1') - (E_2')$ : $(15x - 4x) + (10y - 10y) = 14500 - 8100$
\[ 11x = 6400 \quad\Rightarrow\quad x = \frac{6400}{11} \]

On remplace $x$ dans $(E_1)$ :
\[ 3\times\frac{6400}{11} + 2y = 2900 \;\Rightarrow\; \frac{19200}{11} + 2y = \frac{31900}{11} \;\Rightarrow\; 2y = \frac{12700}{11} \;\Rightarrow\; y = \frac{6350}{11} \]

\textbf{4. Vérification :}
$3x+2y = \frac{19200+12700}{11} = \frac{31900}{11} = 2900$ \checkmark
$2x+5y = \frac{12800+31750}{11} = \frac{44550}{11} = 4050$ \checkmark
Les prix sont positifs, cohérents.

\textbf{5. Conclusion :} Le kilogramme de tomates coûte $\dfrac{6400}{11} \approx 582$ FCFA et le kilogramme d'oignons $\dfrac{6350}{11} \approx 577$ FCFA.
\end{exoresolu}

\begin{methode}[Problèmes « effectif total + autre relation »]
Pour les problèmes du type « il y a $N$ objets/personnes au total, et... » :
\begin{enumerate}
\item La première équation est la \textbf{somme} : $x + y = N$.
\item La deuxième traduit l'autre info (prix total, différence, moyenne pondérée...).
\item La \textbf{substitution} est immédiate : $x = N - y$ (ou $y = N - x$).
\item Vérifier que les solutions sont des \cle{entiers naturels} (on ne vend pas $0,5$ billet).
\end{enumerate}
\end{methode}

\begin{exoresolu}[Billets pour un match]
Pour un match au stade Ahmadou Ahidjo, une école a vendu $120$ billets : « pelouse » à $\SI{500}{FCFA}$ et « tribune » à $\SI{1500}{FCFA}$. Recette totale : $\SI{110000}{FCFA}$.
Combien de billets de chaque catégorie ?
\tcblower
\textbf{Inconnues :} $x$ = nombre de billets pelouse, $y$ = nombre de billets tribune.

\textbf{Système :}
\[ \begin{cases} x + y = 120 \\ 500x + 1500y = 110000 \end{cases} \]
On simplifie la 2ème par $500$ : $x + 3y = 220$.

\textbf{Résolution par substitution :} $x = 120 - y$.
\[ (120 - y) + 3y = 220 \;\Rightarrow\; 120 + 2y = 220 \;\Rightarrow\; 2y = 100 \;\Rightarrow\; y = 50 \]
\[ x = 120 - 50 = 70 \]

\textbf{Vérification :} $70+50=120$ \checkmark ; $70\times500 + 50\times1500 = 35000+75000=110000$ \checkmark. Entiers positifs \checkmark.

\textbf{Conclusion :} L'école a vendu $\mathbf{70}$ billets « pelouse » et $\mathbf{50}$ billets « tribune ».
\end{exoresolu}

\begin{attention}[Les 5 erreurs qui coûtent des points au BEPC]
\begin{enumerate}
\item \textbf{Confondre l'ordre du couple.} Dans $(x\,;y)$, le premier est $x$, le second $y$. Tester $(3\,;2)$ au lieu de $(2\,;3)$ invalide la vérification.
\item \textbf{Oublier de multiplier TOUS les termes.} En combinaison, multiplier une équation par $k$ s'applique aux deux membres : $3x+2y=16 \xrightarrow{\times 2} 6x+4y=32$ (pas $16$).
\item \textbf{Erreur de signe en isolant.} $x + 2y = 7 \Rightarrow x = 7 - 2y$ (pas $+2y$). $2x - y = 1 \Rightarrow y = 2x - 1$ (attention au double changement de signe).
\item \textbf{Ne pas vérifier.} Tester le couple dans \emph{les deux} équations initiales prend 30 secondes et rapporte des points ; cela permet de corriger une erreur de calcul.
\item \textbf{Oublier la définition des inconnues et la conclusion.} « Soit $x$ le prix en FCFA... » est obligatoire. La conclusion doit être une phrase complète avec unités. Sans cela, points de présentation perdus.
\end{enumerate}
\end{attention}

\newpage

\coursec{Exercices}

\courssub{Je vérifie mes connaissances}

\begin{exercice}[QCM : une seule bonne réponse]
Pour chaque question, choisis la bonne réponse et justifie ton choix.
\begin{enumerate}
\item Le couple $(2\,;-1)$ est solution de l'équation :
\begin{enumerate}
\item[a)] $2x + 3y = 7$ \quad b) $3x - 2y = 8$ \quad c) $x - 4y = 5$ \quad d) $5x + y = 8$
\end{enumerate}
\item Le système $\begin{cases} x + y = 5 \\ x - y = 1 \end{cases}$ admet pour solution :
\begin{enumerate}
\item[a)] $(2\,;3)$ \quad b) $(3\,;2)$ \quad c) $(4\,;1)$ \quad d) $(1\,;4)$
\end{enumerate}
\item Graphiquement, un système de deux équations du premier degré à deux inconnues qui n'a \textbf{aucune} solution correspond à :
\begin{enumerate}
\item[a)] deux droites sécantes \quad b) deux droites confondues \quad c) deux droites strictement parallèles \quad d) deux droites perpendiculaires
\end{enumerate}
\item Pour résoudre par combinaison le système $\begin{cases} 3x + 2y = 8 \\ 5x - 2y = 4 \end{cases}$, la première étape la plus simple est :
\begin{enumerate}
\item[a)] multiplier la première ligne par $5$ \quad b) additionner les deux équations membre à membre \quad c) soustraire les deux équations \quad d) exprimer $x$ en fonction de $y$
\end{enumerate}
\end{enumerate}
\end{exercice}

\begin{exercice}[Vrai ou faux, en justifiant]
Réponds par VRAI ou FAUX en justifiant chaque réponse par un calcul ou une propriété du cours.
\begin{enumerate}
\item Le couple $(1\,;2)$ est solution de l'équation $4x - y = 2$.
\item Toute équation du premier degré à deux inconnues $ax + by = c$ (avec $a$ et $b$ non tous les deux nuls) possède une infinité de couples solutions.
\item Le système $\begin{cases} 2x + y = 4 \\ 4x + 2y = 10 \end{cases}$ admet une infinité de solutions.
\item Si deux droites représentant les équations d'un système sont confondues, alors le système n'a pas de solution.
\end{enumerate}
\end{exercice}

\begin{exercice}[Vocabulaire et définitions]
\begin{enumerate}
\item Qu'appelle-t-on \emph{équation du premier degré à deux inconnues} ? Donne un exemple.
\item Qu'appelle-t-on \emph{solution} d'un système de deux équations à deux inconnues ?
\item Cite les deux méthodes algébriques de résolution d'un système de deux équations du premier degré vues en classe.
\item Que représente graphiquement l'ensemble des solutions de l'équation $3x - 2y = 6$ dans un repère du plan ?
\end{enumerate}
\end{exercice}

\begin{exercice}[Calculs directs]
\begin{enumerate}
\item Vérifie si le couple $(3\,;-2)$ est solution de chacune des équations suivantes :
\[ \text{a) } 2x + 5y = -4 \qquad \text{b) } x - 3y = 9 \qquad \text{c) } 4x + y = 14. \]
\item Détermine trois couples solutions de l'équation $x + 2y = 6$.
\item Exprime $y$ en fonction de $x$ dans l'équation $5x - 2y = 10$.
\item Résous mentalement le système $\begin{cases} x + y = 10 \\ x - y = 4 \end{cases}$.
\end{enumerate}
\end{exercice}

\courssub{Je m'entraîne}

\begin{exercice}[Résolution par substitution]
Résous par la méthode de substitution chacun des systèmes suivants, puis vérifie ta solution.
\begin{enumerate}
\item $\begin{cases} y = 2x - 3 \\ 3x + y = 7 \end{cases}$ \qquad 2. $\begin{cases} x - y = 1 \\ 2x + 3y = 12 \end{cases}$ \qquad 3. $\begin{cases} 2x + y = 5 \\ x - 3y = -1 \end{cases}$
\end{enumerate}
\end{exercice}

\begin{exercice}[Résolution par combinaison]
Résous par la méthode de combinaison (élimination) chacun des systèmes suivants, puis vérifie ta solution.
\begin{enumerate}
\item $\begin{cases} 4x + 3y = 11 \\ 2x - 3y = 1 \end{cases}$ \qquad 2. $\begin{cases} 3x + 2y = 12 \\ 5x - 4y = 9 \end{cases}$ \qquad 3. $\begin{cases} 7x - 2y = 16 \\ 3x + 5y = 1 \end{cases}$
\end{enumerate}
\end{exercice}

\begin{exercice}[Approche graphique]
Dans un repère orthonormé $(O\,;\vec{\imath},\vec{\jmath})$ (unité : 1 cm) :
\begin{enumerate}
\item Trace les droites $(d_1)$ et $(d_2)$ d'équations respectives $x + y = 5$ et $2x - y = 1$.
\item Lis graphiquement les coordonnées du point d'intersection de $(d_1)$ et $(d_2)$.
\item Vérifie par le calcul que ce couple est solution du système $\begin{cases} x + y = 5 \\ 2x - y = 1 \end{cases}$.
\item Sur le même repère, trace la droite $(d_3)$ d'équation $x + y = 2$. Que constates-tu pour $(d_1)$ et $(d_3)$ ? Que peux-tu en déduire pour le système $\begin{cases} x + y = 5 \\ x + y = 2 \end{cases}$ ?
\end{enumerate}
\end{exercice}

\begin{exercice}[Nombre de solutions]
Sans résoudre complètement, détermine le nombre de solutions de chacun des systèmes suivants (justifie en comparant les coefficients) :
\begin{enumerate}
\item $\begin{cases} 2x - 3y = 5 \\ 4x - 6y = 10 \end{cases}$ \qquad 2. $\begin{cases} x + 2y = 3 \\ 3x + 6y = 7 \end{cases}$ \qquad 3. $\begin{cases} 2x + y = 4 \\ x - y = 5 \end{cases}$
\end{enumerate}
\end{exercice}

\courssub{Je me prépare au BEPC}

\begin{exercice}[Type BEPC 1 : résolution et discussion]
On considère le système $(S)$ : $\begin{cases} 3x - 2y = 7 \\ 5x + 4y = 19 \end{cases}$.
\begin{enumerate}
\item Résous le système $(S)$ par la méthode de combinaison (élimination), en détaillant toutes les étapes.
\item Vérifie que le couple obtenu est bien solution des deux équations.
\item Reprends la résolution de $(S)$ par la méthode de substitution et montre que l'on obtient la même solution.
\item On considère maintenant le système $(S')$ : $\begin{cases} 2x + ky = 5 \\ 4x + 6y = 10 \end{cases}$ où $k$ est un nombre réel.
\begin{enumerate}
\item Détermine la valeur de $k$ pour laquelle les deux équations de $(S')$ sont proportionnelles.
\item Pour cette valeur de $k$, que peut-on dire du nombre de solutions de $(S')$ ? Donne alors deux couples solutions.
\item Pour $k = 1$, le système $(S')$ admet-il une solution unique ? Justifie sans résoudre.
\end{enumerate}
\end{enumerate}
\end{exercice}

\begin{exercice}[Type BEPC 2 : problème de l'agriculteur]
Un agriculteur de la région de l'Ouest dispose de deux types d'engrais : l'engrais A contient $12\,\%$ d'azote et l'engrais B contient $20\,\%$ d'azote. Il veut obtenir $50$ kg d'un mélange contenant exactement $16\,\%$ d'azote.
\begin{enumerate}
\item On désigne par $x$ la masse (en kg) d'engrais A et par $y$ la masse (en kg) d'engrais B utilisées.
\begin{enumerate}
\item Justifie que $x + y = 50$.
\item Exprime, en fonction de $x$ et $y$, la masse d'azote contenue dans le mélange, puis montre que $0{,}12x + 0{,}20y = 8$.
\end{enumerate}
\item Résous le système obtenu et donne la masse de chaque engrais que l'agriculteur doit utiliser.
\item Vérifie que le mélange obtenu contient bien $16\,\%$ d'azote.
\item Le kilogramme d'engrais A coûte $350$ FCFA et celui d'engrais B coûte $500$ FCFA. Calcule le coût total du mélange.
\end{enumerate}
\end{exercice}

\begin{exercice}[Type BEPC 3 : au marché de Mokolo]
Chez une librairie du marché Mokolo, un carnet et deux stylos coûtent $1\,200$ FCFA ; trois carnets et un stylo coûtent $2\,100$ FCFA.
\begin{enumerate}
\item On note $x$ le prix (en FCFA) d'un carnet et $y$ le prix (en FCFA) d'un stylo. Traduis chaque phrase de l'énoncé par une équation, puis écris le système à résoudre.
\item Résous ce système par la méthode de ton choix et donne le prix d'un carnet et celui d'un stylo.
\item Vérifie ta réponse en calculant le prix de chacun des deux achats de l'énoncé.
\item Amina achète $2$ carnets et $5$ stylos. Elle paie avec un billet de $5\,000$ FCFA. Combien lui rend-on ?
\item Dans un repère orthogonal (1 cm pour $200$ FCFA sur chaque axe), trace les droites d'équations $x + 2y = 1\,200$ et $3x + y = 2\,100$, puis retrouve graphiquement la solution du problème.
\end{enumerate}
\end{exercice}



\newpage

\coursec{Corrigés des exercices}

\begin{corrige}[Exercice 1 — QCM : une seule bonne réponse]
\begin{enumerate}
\item Réponse \textbf{b)}. On remplace $x$ par $2$ et $y$ par $-1$ : $3\times 2 - 2\times(-1) = 6 + 2 = 8$. Vérification des distracteurs : a) $2\times 2+3\times(-1)=4-3=1\neq 7$ ; c) $2-4\times(-1)=2+4=6\neq 5$ ; d) $5\times 2+(-1)=9\neq 8$.
\item Réponse \textbf{b)}. En additionnant membre à membre : $2x=6$, donc $x=3$, puis $y=5-3=2$. Le couple $(3\,;2)$ vérifie bien les deux équations : $3+2=5$ et $3-2=1$.
\item Réponse \textbf{c)}. Deux droites strictement parallèles n'ont aucun point commun : le système n'a aucune solution.
\item Réponse \textbf{b)}. Les coefficients de $y$ sont opposés ($2$ et $-2$) : en additionnant membre à membre, $y$ est éliminé directement et l'on obtient $8x=12$, donc $x=\dfrac{3}{2}$.
\end{enumerate}
\end{corrige}

\begin{corrige}[Exercice 2 — Vrai ou faux, en justifiant]
\begin{enumerate}
\item \textbf{VRAI} : $4\times 1 - 2 = 2$. Le couple $(1\,;2)$ vérifie l'équation.
\item \textbf{VRAI} : comme $a$ et $b$ ne sont pas tous les deux nuls, on peut exprimer une inconnue en fonction de l'autre (par exemple $y=\dfrac{c-ax}{b}$ si $b\neq 0$) ; à chaque valeur de $x$ correspond une valeur de $y$, d'où une infinité de couples solutions (graphiquement, une droite entière).
\item \textbf{FAUX} : en multipliant la première équation par $2$, on obtient $4x+2y=8$, ce qui est incompatible avec $4x+2y=10$. Les droites sont strictement parallèles : le système n'a \textbf{aucune} solution.
\item \textbf{FAUX} : c'est le contraire. Si les droites sont confondues, tous leurs points sont communs : le système admet une \textbf{infinité} de solutions.
\end{enumerate}
\end{corrige}

\begin{corrige}[Exercice 3 — Vocabulaire et définitions]
\begin{enumerate}
\item C'est une égalité de la forme $ax+by=c$, où $a$, $b$, $c$ sont des nombres donnés ($a$ et $b$ non tous les deux nuls) et où $x$ et $y$ sont les inconnues. Exemple : $2x-3y=5$.
\item C'est un couple $(x\,;y)$ de nombres réels qui vérifie \emph{simultanément} les deux équations du système.
\item La méthode de \cle{substitution} et la méthode de \cle{combinaison} (ou élimination par combinaison linéaire).
\item L'ensemble des solutions est représenté par une \cle{droite} : la droite d'équation $3x-2y=6$, c'est-à-dire $y=\dfrac{3}{2}x-3$.
\end{enumerate}
\end{corrige}

\begin{corrige}[Exercice 4 — Calculs directs]
\begin{enumerate}
\item a) $2\times 3 + 5\times(-2) = 6-10 = -4$ : \textbf{oui}. \quad b) $3 - 3\times(-2) = 3+6 = 9$ : \textbf{oui}. \quad c) $4\times 3 + (-2) = 12-2=10 \neq 14$ : \textbf{non}.
\item On choisit des valeurs de $x$ et on calcule $y=\dfrac{6-x}{2}$ : $(0\,;3)$, $(2\,;2)$, $(6\,;0)$ (toute autre réponse correcte est acceptée).
\item $5x-2y=10 \iff 2y = 5x-10 \iff y=\dfrac{5x-10}{2}$, soit $y=\dfrac{5}{2}x-5$.
\item En additionnant membre à membre : $2x=14$, donc $x=7$ et $y=10-7=3$. Solution : $(7\,;3)$.
\end{enumerate}
\end{corrige}

\begin{corrige}[Exercice 5 — Résolution par substitution]
\begin{enumerate}
\item On remplace $y$ par $2x-3$ dans la deuxième équation : $3x+(2x-3)=7$, soit $5x-3=7$, $5x=10$, $x=2$ ; alors $y=2\times 2-3=1$. Vérification : $3\times 2+1=7$. Solution : $(2\,;1)$.
\item La première équation donne $x=y+1$. On substitue : $2(y+1)+3y=12$, soit $5y+2=12$, $5y=10$, $y=2$ ; alors $x=2+1=3$. Vérification : $3-2=1$ et $2\times 3+3\times 2=6+6=12$. Solution : $(3\,;2)$.
\item La première équation donne $y=5-2x$. On substitue : $x-3(5-2x)=-1$, soit $x-15+6x=-1$, $7x=14$, $x=2$ ; alors $y=5-4=1$. Vérification : $2\times 2+1=5$ et $2-3\times 1=-1$. Solution : $(2\,;1)$.
\end{enumerate}
\end{corrige}

\begin{corrige}[Exercice 6 — Résolution par combinaison]
\begin{enumerate}
\item On additionne membre à membre : $(4x+3y)+(2x-3y)=11+1$, soit $6x=12$, donc $x=2$. Alors $4\times 2+3y=11$, soit $3y=3$, $y=1$. Vérification : $8+3=11$ et $4-3=1$. Solution : $(2\,;1)$.
\item On multiplie la première équation par $2$ : $6x+4y=24$. On additionne avec la deuxième : $11x=33$, donc $x=3$. Alors $3\times 3+2y=12$, soit $2y=3$, $y=\dfrac{3}{2}$. Vérification : $9+3=12$ et $15-4\times\dfrac{3}{2}=15-6=9$. Solution : $\left(3\,;\dfrac{3}{2}\right)$.
\item On multiplie la première équation par $5$ et la deuxième par $2$ : $35x-10y=80$ et $6x+10y=2$. On additionne : $41x=82$, donc $x=2$. Alors $7\times 2-2y=16$, soit $-2y=2$, $y=-1$. Vérification : $14-2\times(-1)=16$ et $3\times 2+5\times(-1)=6-5=1$. Solution : $(2\,;-1)$.
\end{enumerate}
\end{corrige}

\begin{corrige}[Exercice 7 — Approche graphique]
\begin{enumerate}
\item $(d_1)$ : $y=-x+5$, passe par $(0\,;5)$ et $(5\,;0)$. $(d_2)$ : $y=2x-1$, passe par $(0\,;-1)$ et $(2\,;3)$. On place ces points et on trace les deux droites à la règle.
\item Le point d'intersection lu sur le graphique a pour coordonnées $(2\,;3)$.
\item Vérification : $2+3=5$ et $2\times 2-3=1$. Le couple $(2\,;3)$ est bien la solution du système.
\item $(d_3)$ : $y=-x+2$, passe par $(0\,;2)$ et $(2\,;0)$. Les droites $(d_1)$ et $(d_3)$ ont le même coefficient directeur $-1$ et des ordonnées à l'origine différentes ($5$ et $2$) : elles sont \textbf{strictement parallèles}. Le système n'a donc \textbf{aucune solution} (on ne peut pas avoir à la fois $x+y=5$ et $x+y=2$).
\end{enumerate}
\end{corrige}

\begin{corrige}[Exercice 8 — Nombre de solutions]
\begin{enumerate}
\item La deuxième équation est le double de la première ($4=2\times 2$, $-6=2\times(-3)$, $10=2\times 5$) : les deux équations sont équivalentes, les droites sont \textbf{confondues}. Le système admet une \textbf{infinité de solutions}.
\item On a $3\times 1=3$ et $3\times 2=6$, mais $3\times 3=9\neq 7$ : les coefficients de $x$ et $y$ sont proportionnels, mais pas les seconds membres. Les droites sont \textbf{strictement parallèles} : \textbf{aucune solution}.
\item On a $\dfrac{2}{1}\neq \dfrac{1}{-1}$ : les coefficients ne sont pas proportionnels, les droites sont \textbf{sécantes}. Le système admet une \textbf{solution unique} (c'est $(3\,;-2)$, ce qu'on peut vérifier : $6-2=4$ et $3+2=5$).
\end{enumerate}
\end{corrige}

\begin{corrige}[Exercice 9 — Type BEPC 1 : résolution et discussion]
\begin{enumerate}
\item On multiplie la première équation par $2$ : $6x-4y=14$. On additionne membre à membre avec la deuxième équation : $6x+5x=14+19$, soit $11x=33$, donc $x=3$. On reporte dans la première équation : $3\times 3-2y=7$, soit $-2y=-2$, $y=1$. Solution : $(3\,;1)$.
\item Vérification : $3\times 3-2\times 1=9-2=7$ et $5\times 3+4\times 1=15+4=19$. Le couple $(3\,;1)$ vérifie les deux équations.
\item La première équation donne $2y=3x-7$, soit $y=\dfrac{3x-7}{2}$. On substitue dans la deuxième : $5x+4\times\dfrac{3x-7}{2}=19$, soit $5x+2(3x-7)=19$, $5x+6x-14=19$, $11x=33$, $x=3$ ; alors $y=\dfrac{9-7}{2}=1$. On retrouve bien la solution $(3\,;1)$.
\item \begin{enumerate}
\item Les deux équations sont proportionnelles lorsque $4=2\times 2$, $6=2k$ et $10=2\times 5$, ce qui donne $k=3$.
\item Pour $k=3$, la deuxième équation est le double de la première : les droites sont confondues, le système admet une \textbf{infinité de solutions}. L'équation $2x+3y=5$ donne par exemple les couples $(1\,;1)$ et $(4\,;-1)$.
\item Pour $k=1$, on compare : $\dfrac{2}{4}=\dfrac{1}{2}$ et $\dfrac{k}{6}=\dfrac{1}{6}$ ; comme $\dfrac{1}{2}\neq\dfrac{1}{6}$, les coefficients ne sont pas proportionnels : les droites sont sécantes, le système admet une \textbf{solution unique}.
\end{enumerate}
\end{enumerate}
\textbf{Barème indicatif (sur 5 pts)} : 1 pt pour la résolution par combinaison ; $0{,}5$ pt pour la vérification ; 1 pt pour la substitution ; $0{,}75$ pt pour $k=3$ ; $0{,}75$ pt pour l'infinité de solutions et les deux couples ; 1 pt pour la justification du cas $k=1$.\par
\textbf{Points de vigilance} : exiger l'écriture de l'opération effectuée sur les lignes (« on multiplie la première équation par 2 ») ; la vérification doit porter sur les \emph{deux} équations ; pour la discussion, comparer explicitement les rapports de coefficients avant de conclure.
\end{corrige}

\begin{corrige}[Exercice 10 — Type BEPC 2 : problème de l'agriculteur]
\begin{enumerate}
\item \begin{enumerate}
\item La masse totale du mélange est la somme des masses des deux engrais, et elle doit valoir $50$ kg : $x+y=50$.
\item L'engrais A apporte $0{,}12x$ kg d'azote et l'engrais B apporte $0{,}20y$ kg. Le mélange doit contenir $16\,\%$ de $50$ kg d'azote, soit $0{,}16\times 50=8$ kg. D'où $0{,}12x+0{,}20y=8$.
\end{enumerate}
\item On résout $\begin{cases} x+y=50 \\ 0{,}12x+0{,}20y=8 \end{cases}$. De la première équation : $x=50-y$. On substitue : $0{,}12(50-y)+0{,}20y=8$, soit $6-0{,}12y+0{,}20y=8$, $0{,}08y=2$, $y=\dfrac{2}{0{,}08}=25$ ; alors $x=50-25=25$. L'agriculteur doit utiliser \textbf{25 kg d'engrais A et 25 kg d'engrais B}.
\item Masse d'azote : $0{,}12\times 25+0{,}20\times 25=3+5=8$ kg, et $\dfrac{8}{50}=0{,}16=16\,\%$. Le mélange est conforme.
\item Coût total : $25\times 350+25\times 500=8\,750+12\,500=21\,250$ FCFA.
\end{enumerate}
\textbf{Barème indicatif (sur 5 pts)} : 1 pt pour la modélisation (les deux équations) ; 2 pts pour la résolution détaillée ; 1 pt pour la vérification du pourcentage ; 1 pt pour le coût total.\par
\textbf{Points de vigilance} : bien convertir les pourcentages en coefficients décimaux ($12\,\%=0{,}12$) ; calculer d'abord la masse d'azote attendue ($8$ kg) ; conclure par une phrase répondant à la question posée, avec les unités.
\end{corrige}

\begin{corrige}[Exercice 11 — Type BEPC 3 : au marché de Mokolo]
\begin{enumerate}
\item « Un carnet et deux stylos coûtent $1\,200$ FCFA » se traduit par $x+2y=1\,200$ ; « trois carnets et un stylo coûtent $2\,100$ FCFA » se traduit par $3x+y=2\,100$. Le système est : $\begin{cases} x+2y=1\,200 \\ 3x+y=2\,100 \end{cases}$.
\item Par substitution : la première équation donne $x=1\,200-2y$. On reporte : $3(1\,200-2y)+y=2\,100$, soit $3\,600-6y+y=2\,100$, $-5y=-1\,500$, $y=300$ ; alors $x=1\,200-2\times 300=600$. Un carnet coûte \textbf{600 FCFA} et un stylo \textbf{300 FCFA}.
\item Vérification : $600+2\times 300=1\,200$ FCFA et $3\times 600+300=2\,100$ FCFA. Les deux achats correspondent bien à l'énoncé.
\item Amina paie $2\times 600+5\times 300=1\,200+1\,500=2\,700$ FCFA. On lui rend $5\,000-2\,700=2\,300$ FCFA.
\item La droite $x+2y=1\,200$ passe par $(0\,;600)$ et $(1\,200\,;0)$ ; la droite $3x+y=2\,100$ passe par $(0\,;2\,100)$ et $(700\,;0)$. Avec l'échelle 1 cm pour 200 FCFA : $600 \rightarrow 3$ cm, $1\,200 \rightarrow 6$ cm, $2\,100 \rightarrow 10{,}5$ cm, $700 \rightarrow 3{,}5$ cm. Les deux droites se coupent au point de coordonnées $(600\,;300)$, ce qui confirme graphiquement la solution trouvée par le calcul.
\end{enumerate}
\textbf{Barème indicatif (sur 6 pts)} : 1 pt pour la traduction en système ; 2 pts pour la résolution ; $0{,}5$ pt pour la vérification ; 1 pt pour la monnaie rendue ; $1{,}5$ pt pour le graphique (points bien placés, droites tracées, intersection lue).\par
\textbf{Points de vigilance} : définir clairement les inconnues avant d'écrire le système ; respecter l'échelle imposée sur les deux axes ; la lecture graphique doit être confirmée par le calcul, pas l'inverse.
\end{corrige}

\newpage

\coursec{Fiche bilan}

\begin{popfigure}
\begin{tikzpicture}[mindmap, grow cyclic, every node/.style={concept, font=\small},
  root concept/.append style={concept color=popblue, font=\bfseries},
  level 1/.append style={level distance=3.6cm, sibling angle=51.4},
  level 1 concept/.append style={font=\footnotesize}]
\node[root concept]{Équations du 1\textsuperscript{er} degré dans $\mathbb{R}\times\mathbb{R}$}
  child[concept color=poporange]{ node {Équation $ax+by=c$ : inconnues $x,y$ ; solution = couple $(x\,;y)$ ; graphe = droite} }
  child[concept color=popgreen]{ node[align=center]{Système $\begin{cases}ax+by=c\\a'x+b'y=c'\end{cases}$ : solution = couple vérifiant les DEUX équations} }
  child[concept color=poppurple]{ node {Approche graphique : tracer 2 droites, lire le point d'intersection} }
  child[concept color=poppink]{ node {Substitution : isoler une inconnue, remplacer dans l'autre équation} }
  child[concept color=popgold]{ node {Élimination : multiplier, additionner pour supprimer une inconnue} }
  child[concept color=popblueL]{ node {Nombre de solutions : $ab'-a'b\neq 0$ unique ; $=0$ aucune ou infinité} }
  child[concept color=popgreenL]{ node {Problèmes concrets : choisir les inconnues, traduire, résoudre, vérifier, conclure} };
\end{tikzpicture}
\legende{Carte mentale de la leçon : équations et systèmes du premier degré à deux inconnues.}
\end{popfigure}

\begin{aretenir}[Bilan : définitions clés]
\begin{itemize}
  \item Une \cle{équation du premier degré à deux inconnues} est une équation qui peut s'écrire sous la forme $ax+by=c$, où $a$, $b$, $c$ sont des réels donnés avec $(a\,;b)\neq(0\,;0)$, et où $x$ et $y$ sont les inconnues.
  \item Une \cle{solution} d'une telle équation est un \cle{couple} $(x\,;y)$ de réels qui rend l'égalité vraie. L'ordre compte : en général $(2\,;5)\neq(5\,;2)$.
  \item L'ensemble des solutions de $ax+by=c$ est représenté graphiquement par une \cle{droite} du plan muni d'un repère.
  \item Un \cle{système de deux équations} du premier degré à deux inconnues s'écrit : $\begin{cases} ax+by=c \\ a'x+b'y=c' \end{cases}$. Le résoudre, c'est trouver tous les couples $(x\,;y)$ solutions \textbf{communes} aux deux équations.
  \item Le réel $\Delta = ab'-a'b$ est le \cle{déterminant} du système. Il permet de prévoir le nombre de solutions sans résoudre entièrement.
\end{itemize}
\end{aretenir}

\begin{aretenir}[Bilan : formules et résultats à connaître par c\oe ur]
\begin{center}
\begin{tabularx}{\linewidth}{|l|X|}
\hline
\rowcolor{popblueL} \textbf{Situation} & \textbf{Conclusion} \\
\hline
$\Delta = ab'-a'b \neq 0$ & Le système admet une \textbf{unique solution} (les deux droites sont sécantes). \\
\hline
$\Delta = 0$ et les coefficients sont proportionnels, $c$ et $c'$ compris (équations équivalentes) & \textbf{Infinité de solutions} (droites confondues). \\
\hline
$\Delta = 0$ sans proportionnalité complète & \textbf{Aucune solution} (droites strictement parallèles). \\
\hline
Solution unique ($\Delta\neq 0$) & $x=\dfrac{cb'-c'b}{ab'-a'b}$ et $y=\dfrac{ac'-a'c}{ab'-a'b}$ \\
\hline
\end{tabularx}
\end{center}
\emph{Remarque :} au BEPC, on retrouve le plus souvent la solution par substitution ou par élimination plutôt qu'en appliquant directement les formules ; ces formules servent surtout à vérifier.
\end{aretenir}

\begin{methode}[Bilan : méthodes express]
\begin{itemize}
  \item \textbf{Substitution} (idéale si un coefficient vaut $1$ ou $-1$) :
  \begin{enumerate}
    \item isoler une inconnue dans l'une des équations ;
    \item la remplacer dans l'autre équation ;
    \item résoudre l'équation à une seule inconnue obtenue ;
    \item calculer la valeur de l'autre inconnue ;
    \item vérifier le couple trouvé dans les deux équations de départ.
  \end{enumerate}
  \item \textbf{Élimination (combinaison linéaire)} :
  \begin{enumerate}
    \item multiplier les équations par des nombres bien choisis pour obtenir des coefficients opposés devant une même inconnue ;
    \item additionner membre à membre pour éliminer cette inconnue ;
    \item résoudre l'équation à une inconnue obtenue ;
    \item remonter à l'autre inconnue en remplaçant dans une équation de départ ;
    \item vérifier.
  \end{enumerate}
  \item \textbf{Méthode graphique} : écrire chaque équation sous la forme $y=\ldots$, tracer les deux droites (deux points suffisent pour chacune), puis lire les coordonnées du point d'intersection : c'est la solution du système.
  \item \textbf{Problème concret} : nommer les inconnues $\to$ traduire chaque information de l'énoncé en équation $\to$ résoudre le système $\to$ \textbf{vérifier} la solution dans l'énoncé $\to$ conclure par une phrase complète avec les unités.
\end{itemize}
\end{methode}

\begin{attention}[Pièges fréquents au BEPC]
\begin{itemize}
  \item Ne pas confondre l'ordre des coordonnées : la solution s'écrit $(x\,;y)$, $x$ d'abord.
  \item Dans la substitution, remplacer dans \textbf{l'autre} équation (sinon on obtient $0=0$, qui ne donne aucune information).
  \item Dans l'élimination, multiplier \textbf{tous} les termes de l'équation, y compris le second membre.
  \item Attention aux signes lors de l'addition membre à membre : $-(-3y)=+3y$.
  \item Une solution trouvée graphiquement doit être confirmée par le calcul si l'énoncé demande une valeur exacte.
  \item Toujours vérifier le couple solution dans les \textbf{deux} équations de départ.
\end{itemize}
\end{attention}

\begin{aretenir}[Checklist avant le BEPC]
\begin{itemize}
  \item[$\square$] Je sais reconnaître une équation du premier degré à deux inconnues $ax+by=c$.
  \item[$\square$] Je sais vérifier si un couple $(x\,;y)$ est solution d'une équation ou d'un système.
  \item[$\square$] Je sais qu'une solution est un \textbf{couple ordonné} et je respecte l'ordre $(x\,;y)$.
  \item[$\square$] Je sais résoudre un système par \textbf{substitution} sans erreur de signe.
  \item[$\square$] Je sais résoudre un système par \textbf{élimination} (combinaison linéaire).
  \item[$\square$] Je sais interpréter graphiquement : point d'intersection = solution du système.
  \item[$\square$] Je sais calculer le déterminant $\Delta=ab'-a'b$ et conclure sur le nombre de solutions.
  \item[$\square$] Je n'oublie jamais de \textbf{vérifier} ma solution dans les DEUX équations de départ.
  \item[$\square$] Je sais traduire un énoncé concret (prix, âges, parts, transports) en système de deux équations.
  \item[$\square$] Je rédige une conclusion claire, avec une phrase complète et les unités.
\end{itemize}
\end{aretenir}

\findecours

\end{document}
